% =====================================================================
%  Capítulo 16 · Biología de sistemas y redes
% =====================================================================
\chapter{Biología de sistemas y redes}\label{cap:16}

\epigrafe{Este enfoque puede revelar los bloques de construcción básicos de la mayoría de las redes.}{\textcite{c16-milo2002}}

\begin{objetivos}
\begin{itemize}
  \item Representar un interactoma como un grafo y calcular, a mano y con NetworkX, grado, coeficiente de agrupamiento, caminos mínimos, intermediación y modularidad.
  \item Explicar cómo se obtienen las interacciones proteína-proteína (doble híbrido, purificación por afinidad con espectrometría de masas), qué errores introduce cada método y cómo los integran STRING y BioGRID.
  \item Derivar las propiedades de los modelos de Erdős-Rényi, Watts-Strogatz y Barabási-Albert, y evaluar con espíritu crítico la afirmación de que las redes biológicas son ``libres de escala''.
  \item Construir modelos de ecuaciones diferenciales de la expresión génica con funciones de Hill, y deducir de ellos el tiempo de respuesta, el efecto de la autorregulación negativa y la función del \ingles{feed-forward loop}.
  \item Analizar la estabilidad de un circuito (nulclinas, jacobiano, autovalores) para explicar la biestabilidad del interruptor genético y las oscilaciones del \ingles{repressilator}.
  \item Simular el ruido de la expresión con el algoritmo de Gillespie e inferir redes de regulación a partir de datos de expresión con GENIE3.
\end{itemize}
\end{objetivos}

Hasta aquí, casi todo el libro ha tratado a los genes y a las proteínas como entidades individuales: una secuencia que se alinea, una estructura que se predice, un transcrito cuya abundancia cambia entre dos condiciones (\cref{cap:11}). Pero ninguna proteína trabaja sola. La quinasa FUS3 de la levadura, por ejemplo, solo transmite la señal de la feromona sexual porque una proteína andamio la acerca a su activadora, y solo produce una respuesta porque fosforila a un factor de transcripción que, a su vez, enciende decenas de genes. La función biológica es una propiedad de las \emph{relaciones} entre moléculas, y el lenguaje natural para describir relaciones es el de los grafos.

La biología de sistemas estudia esas relaciones con dos miradas complementarias, que corresponden a las dos lecciones de este capítulo. La primera es \emph{estructural}: ¿cómo está cableada la célula?, ¿qué proteínas son centrales?, ¿se agrupan en módulos? La abordaremos con las redes de interacción proteína-proteína, sus métodos experimentales, sus bases de datos y la teoría de grafos que permite medirlas (\cref{sec:16-ppi}). La segunda es \emph{dinámica}: dado un pequeño circuito de genes que se regulan entre sí, ¿qué hace en el tiempo? Responderemos con ecuaciones diferenciales, análisis de estabilidad y simulación estocástica, y terminaremos preguntándonos cómo reconstruir el circuito a partir de datos de expresión (\cref{sec:16-grn}). Ambas miradas comparten una idea que \textcite{c16-barabasi2004} resumieron bien: entender la organización funcional de la célula exige pasar de catálogos de componentes a mapas de sus interacciones.

% ---------------------------------------------------------------------
\section{Redes de interacción proteína-proteína}\label{sec:16-ppi}
% ---------------------------------------------------------------------
\enlacerepo{16.1 · Redes proteína-proteína}

\begin{intuicion}[El mapa del metro]
Un mapa del metro no respeta distancias ni la forma de las calles: solo dice qué estaciones están conectadas. Y sin embargo, con esa información mínima podemos responder preguntas muy útiles. ¿Cuántos transbordos hay entre dos estaciones? ¿Qué estación, si se cierra, desconecta media ciudad? ¿Qué líneas forman barrios bien comunicados entre sí pero mal comunicados con el resto? Un mapa de interacciones proteicas funciona igual: olvida la química, la estructura tridimensional y la cinética, y guarda solo quién toca a quién. A cambio, nos permite preguntar qué proteínas son ``estaciones de transbordo'', cuán lejos está una proteína de otra o qué grupos de proteínas trabajan juntos.
\end{intuicion}

\subsection{El vocabulario de los grafos}

\begin{definicion}{Grafo y matriz de adyacencia}{16-grafo}
Un \emph{grafo} $G=(V,E)$ es un conjunto de $n=|V|$ nodos (vértices) y $m=|E|$ aristas (enlaces), cada una un par de nodos. Si los pares no tienen orden, el grafo es \emph{no dirigido} (interacciones físicas); si lo tienen, es \emph{dirigido} (regulación: el factor de transcripción $i$ actúa sobre el gen $j$). El grafo se codifica en la matriz de adyacencia $A\in\{0,1\}^{n\times n}$:
\begin{equation}\label{eq:16-adj}
A_{ij}=\begin{cases}1 & \text{si hay una arista entre } i \text{ y } j,\\ 0 & \text{en otro caso.}\end{cases}
\end{equation}
En un grafo no dirigido $A$ es simétrica; en uno ponderado, $A_{ij}$ puede ser un peso real (por ejemplo, una puntuación de confianza).
\end{definicion}

\begin{simbolos}
\simbolo{$V,\ E$}{Conjunto de nodos (proteínas) y de aristas (interacciones).}
\simbolo{$n,\ m$}{Número de nodos y de aristas.}
\simbolo{$A_{ij}$}{Entrada de la matriz de adyacencia: 1 si $i$ y $j$ interactúan.}
\end{simbolos}

La cantidad más sencilla que podemos asociar a un nodo es su \emph{grado}, el número de vecinos:
\begin{equation}\label{eq:16-grado}
k_i=\sum_{j=1}^{n}A_{ij},\qquad \sum_{i=1}^{n}k_i = 2m,\qquad \langle k\rangle=\frac{2m}{n}.
\end{equation}
\begin{simbolos}
\simbolo{$k_i$}{Grado del nodo $i$ (número de interacciones de la proteína).}
\simbolo{$\langle k\rangle$}{Grado medio de la red.}
\end{simbolos}
La segunda igualdad (el ``lema del apretón de manos'') se debe a que cada arista aporta un vecino a cada uno de sus dos extremos. La fracción de nodos con grado $k$, $P(k)$, es la \emph{distribución de grado}, y será protagonista más adelante.

La matriz de adyacencia no es solo una tabla: sus potencias cuentan caminos. El elemento $(A^2)_{ij}=\sum_{l}A_{il}A_{lj}$ suma 1 por cada nodo $l$ vecino a la vez de $i$ y de $j$, es decir, cuenta los recorridos de longitud 2 entre ellos. Por inducción,
\begin{equation}\label{eq:16-caminos}
(A^{\ell})_{ij} = \text{número de recorridos de longitud } \ell \text{ entre } i \text{ y } j.
\end{equation}
En particular, $(A^3)_{ii}$ cuenta los recorridos que salen de $i$ y vuelven a él en tres pasos; cada triángulo que contiene a $i$ se puede recorrer en dos sentidos, así que el número de triángulos de $i$ es $t_i=(A^3)_{ii}/2$. Esto nos lleva a una medida de lo ``apiñado'' que está el vecindario de un nodo.

\begin{definicion}{Coeficiente de agrupamiento}{16-clustering}
El coeficiente de agrupamiento local de un nodo con $k_i\ge2$ es la fracción de pares de sus vecinos que también son vecinos entre sí:
\begin{equation}\label{eq:16-clustering}
C_i=\frac{t_i}{\binom{k_i}{2}}=\frac{2\,t_i}{k_i(k_i-1)},\qquad C=\frac{1}{n}\sum_i C_i .
\end{equation}
\end{definicion}
\begin{simbolos}
\simbolo{$t_i$}{Número de triángulos que contienen al nodo $i$, igual a $(A^3)_{ii}/2$.}
\simbolo{$\binom{k_i}{2}$}{Número de pares posibles entre los $k_i$ vecinos de $i$.}
\simbolo{$C$}{Coeficiente de agrupamiento medio de la red (con $C_i=0$ si $k_i<2$).}
\end{simbolos}

La \emph{distancia} $d_{ij}$ entre dos nodos es la longitud del camino más corto que los une; se calcula por búsqueda en anchura en tiempo $O(n+m)$ desde cada origen. Su promedio sobre todos los pares conectados, $L$, y su máximo, el \emph{diámetro}, resumen cuán ``compacta'' es la red. Finalmente, la \emph{intermediación} (\ingles{betweenness}) de un nodo mide cuántos caminos mínimos pasan por él:
\begin{equation}\label{eq:16-betweenness}
b_v=\sum_{s\neq v\neq t}\frac{\sigma_{st}(v)}{\sigma_{st}},
\end{equation}
\begin{simbolos}
\simbolo{$\sigma_{st}$}{Número de caminos mínimos entre $s$ y $t$.}
\simbolo{$\sigma_{st}(v)$}{Cuántos de ellos pasan por $v$.}
\simbolo{$b_v$}{Intermediación: nodos ``puente'' por los que fluye mucha comunicación.}
\end{simbolos}
Un nodo puede tener grado bajo y aun así intermediación altísima, si es el único enlace entre dos regiones de la red. El grado mide popularidad local; la intermediación, importancia estratégica global.

\begin{figure}[t]
\centering
\begin{tikzpicture}[font=\sffamily\small,
  v/.style={circle,draw=white,thick,fill=#1,minimum size=7mm,inner sep=0pt,text=white,font=\sffamily\bfseries\small}]
  \node[v=azul] (A) at (0,1.5) {A};
  \node[v=azul] (B) at (0,0) {B};
  \node[v=naranja] (C) at (1.3,0.75) {C};
  \node[v=naranja] (D) at (2.7,0.75) {D};
  \node[v=aqua] (E) at (3.8,1.6) {E};
  \node[v=naranja] (F) at (3.8,-0.1) {F};
  \node[v=aqua] (G) at (5.2,-0.1) {G};
  \begin{scope}[on background layer]
  \foreach \a/\b in {A/B,A/C,B/C,C/D,D/E,D/F,E/F,F/G}{\draw[tinta2,thick] (\a)--(\b);}
  \fill[amarillo!25] (A.center)--(B.center)--(C.center)--cycle;
  \end{scope}
  \node[etiqueta,text=tinta2,anchor=north,align=center] at (2.6,-0.75) {naranja: intermediación $b_v>0$\\ amarillo: el triángulo A--B--C};
  \matrix[matrix of math nodes, anchor=west, nodes={minimum width=5.2mm, minimum height=4.6mm, font=\small, inner sep=0pt},
     left delimiter={[}, right delimiter={]}] (M) at (7.0,0.75) {
    0&1&1&0&0&0&0\\
    1&0&1&0&0&0&0\\
    1&1&0&1&0&0&0\\
    0&0&1&0&1&1&0\\
    0&0&0&1&0&1&0\\
    0&0&0&1&1&0&1\\
    0&0&0&0&0&1&0\\};
  \foreach \l [count=\i] in {A,B,C,D,E,F,G}{
    \node[etiqueta,text=gris] at ($(M-\i-1.west)+(-0.45,0)$) {\l};
    \node[etiqueta,text=gris,above=2pt] at (M-1-\i.north) {\l};}
  \begin{scope}[on background layer]
  \fill[amarillo!30,rounded corners=1pt] (M-3-1.north west) rectangle (M-3-7.south east);
  \end{scope}
  \node[etiqueta,text=tinta2,anchor=west] at ($(M-3-7.east)+(0.35,0)$) {$k_C=3$};
\end{tikzpicture}
\caption[Un grafo pequeño y su matriz de adyacencia]{Un grafo de siete proteínas y su matriz de adyacencia (\cref{eq:16-adj}). La fila resaltada suma el grado de C, $k_C=3$. Los nodos A, B y E tienen $C_i=1$ (sus dos vecinos están conectados entre sí), mientras que C, D y F tienen $C_i=1/3$. Aunque C, D y F tienen el mismo grado, D tiene la mayor intermediación (9 caminos mínimos entre pares la atraviesan), porque es el único puente entre los dos grupos. Cortar la arista C--D da la mejor partición en comunidades (\cref{ej:16-modularidad}).}
\label{fig:16-grafo}
\end{figure}

\begin{ejemplo}{Medir a mano un grafo pequeño}{16-toy}
En el grafo de la \cref{fig:16-grafo} hay $n=7$ nodos y $m=8$ aristas. Los grados son $(2,2,3,3,2,3,1)$ para A,\dots,G; suman $16=2m$, y $\langle k\rangle=16/7\approx2{,}29$.
\begin{itemize}
  \item \textbf{Caminos.} $(A^2)_{\mathrm{A,D}}=1$: hay un único recorrido de longitud 2 de A a D (por C). $(A^3)_{\mathrm{A,E}}=1$: el recorrido A--C--D--E.
  \item \textbf{Agrupamiento.} C tiene vecinos A, B y D; de los tres pares posibles solo A--B está unido, así que $C_{\mathrm{C}}=1/3$. Promediando, $C=(1+1+\tfrac13+\tfrac13+1+\tfrac13+0)/7=4/7\approx0{,}571$.
  \item \textbf{Distancias.} La distancia media entre los 21 pares es $L\approx2{,}05$ y el diámetro es 4 (de A o B a G).
  \item \textbf{Intermediación.} $b_{\mathrm{C}}=8$, $b_{\mathrm{D}}=9$, $b_{\mathrm{F}}=5$, y cero para los demás. D gana porque todo camino entre $\{$A, B, C$\}$ y $\{$E, F, G$\}$ la cruza.
\end{itemize}
Todas estas cifras se reproducen con \archivo{figuras/cap16/generar.py}.
\end{ejemplo}

\subsection{Cómo se mide un interactoma}

Los grafos de interacción no se observan directamente: se \emph{construyen} a partir de experimentos, y cada experimento deja su huella en la red resultante. Dos familias de métodos dominaron la primera década de la interactómica.

\paragraph{Doble híbrido en levadura (Y2H)} Un factor de transcripción como Gal4 tiene dos dominios separables: uno que se une al ADN (DBD) y otro que activa la transcripción (AD). Por separado no hacen nada; si se acercan, reconstituyen un activador funcional. El método fusiona una proteína ``cebo'' al DBD y otra ``presa'' al AD: si cebo y presa interactúan, el activador se reconstituye en el promotor y se expresa un gen reportero que permite a la levadura crecer en un medio selectivo. \textcite{c16-uetz2000} aplicaron el método a escala de genoma con dos estrategias (una matriz ordenada de unos 6000 transformantes y un cribado de bibliotecas mezcladas) y detectaron 957 interacciones putativas entre \num{1004} proteínas de \especie{Saccharomyces cerevisiae}. Casi al mismo tiempo, \textcite{c16-ito2001} completaron un cribado exhaustivo que identificó \num{4549} interacciones entre \num{3278} proteínas y notaron algo inquietante: sus datos \emph{no se solapaban en gran medida} con los de Uetz y colaboradores.

\paragraph{Purificación por afinidad y espectrometría de masas (AP-MS)} En lugar de pares, AP-MS mide complejos. Se marca una proteína cebo con una etiqueta de afinidad, se purifica en condiciones suaves junto con todo lo que arrastra, y los acompañantes se identifican por espectrometría de masas. \textcite{c16-gavin2002} usaron purificación en tándem (TAP) con \num{1739} genes de levadura, purificaron 589 ensamblajes proteicos y definieron 232 complejos multiproteicos distintos. El resultado no es una lista de pares sino de conjuntos, y convertirlos en aristas exige una decisión de modelado (\cref{fig:16-metodos}): en el modelo \emph{estrella} (\ingles{spoke}) se une el cebo a cada presa; en el modelo \emph{matriz} (\ingles{matrix}) se unen todos con todos. Con un cebo y cinco presas, el primero produce 5 aristas y el segundo $\binom{6}{2}=15$, la mayoría no verificadas como contactos directos.

\begin{figure}[t]
\centering
\begin{tikzpicture}[font=\sffamily\small,
   prot/.style={draw=#1!70!black,fill=#1!25,thick,rounded corners=3pt,minimum height=6mm,inner sep=2pt},
   p/.style={circle,fill=#1,minimum size=4.5mm,inner sep=0pt,text=white,font=\sffamily\scriptsize\bfseries}]
  % ---- Y2H
  \node[etiqueta,font=\sffamily\bfseries\small] at (2.2,3.3) {Doble híbrido (Y2H)};
  \draw[very thick,tinta2] (-0.3,0) -- (4.9,0);
  \fill[amarillo!50] (0.2,-0.12) rectangle (1.4,0.12);
  \node[etiqueta,below] at (0.8,-0.15) {UAS};
  \draw[very thick,verde,-{Stealth}] (2.3,0) -- (2.3,0.45) -- (3.0,0.45);
  \fill[verde!30] (2.3,-0.12) rectangle (4.6,0.12);
  \node[etiqueta,below] at (3.45,-0.15) {gen reportero};
  \node[prot=azul,minimum width=1.2cm] (dbd) at (0.8,0.48) {DBD};
  \node[prot=azul,minimum width=1.0cm,anchor=south] (bait) at (dbd.north) {cebo};
  \node[prot=naranja,minimum width=1.0cm,anchor=west] (prey) at (bait.east) {presa};
  \node[prot=naranja,minimum width=1.0cm,anchor=south] (ad) at (prey.north) {AD};
  \draw[flecha=aqua] (ad.east) to[bend left=25] node[right,etiqueta,text=aqua!60!black,align=left,pos=0.4]{recluta la\\maquinaria} (3.2,0.6);
  % ---- AP-MS
  \begin{scope}[xshift=6.9cm]
  \node[etiqueta,font=\sffamily\bfseries\small] at (2.5,3.3) {AP-MS: complejo purificado};
  \node[p=azul] (c) at (0.8,1.2) {C};
  \foreach \i/\a in {1/90,2/18,3/-54,4/-126,5/162}{\node[p=naranja] (q\i) at ($(c)+(\a:0.5)$) {};}
  \draw[amarillo!70!black,thick] (q4.south east) -- ++(0.1,-0.45) node[below,etiqueta]{etiqueta TAP};
  \node[etiqueta,text=tinta2] at (0.8,2.2) {cebo + 5 presas};
  % spoke
  \node[p=azul] (s) at (3.0,1.9) {C};
  \foreach \i/\a in {1/90,2/18,3/-54,4/-126,5/162}{\node[p=naranja,minimum size=3mm] (s\i) at ($(s)+(\a:0.6)$) {};
     \draw[azul,thick] (s) -- (s\i);}
  \node[etiqueta,text=tinta2,anchor=west,align=left] at (3.75,1.9) {estrella:\\5 aristas};
  % matrix
  \node[p=azul] (t) at (3.0,0.15) {C};
  \foreach \i/\a in {1/90,2/18,3/-54,4/-126,5/162}{\node[p=naranja,minimum size=3mm] (t\i) at ($(t)+(\a:0.6)$) {};}
  \begin{scope}[on background layer]
  \foreach \x in {t,t1,t2,t3,t4,t5}{\foreach \y in {t,t1,t2,t3,t4,t5}{\draw[naranja!60,thin] (\x) -- (\y);}}
  \end{scope}
  \node[etiqueta,text=tinta2,anchor=west,align=left] at (3.75,0.15) {matriz:\\15 aristas};
  \end{scope}
\end{tikzpicture}
\caption[Doble híbrido y AP-MS]{Los dos grandes métodos experimentales de la interactómica. \textbf{Izquierda}: en el doble híbrido, el cebo fusionado al dominio de unión al ADN (DBD) y la presa fusionada al dominio de activación (AD) reconstituyen un activador solo si interactúan; la lectura es un par binario. \textbf{Derecha}: AP-MS purifica un complejo completo a partir de un cebo etiquetado. Traducirlo a aristas exige elegir un modelo: \emph{estrella} (subestima contactos entre presas) o \emph{matriz} (sobreestima: la mayoría de los 15 pares no se tocan físicamente). Esa elección explica en parte por qué los complejos grandes aparecen en las redes como densas ``bolas'' de aristas.}
\label{fig:16-metodos}
\end{figure}

\subsection{Falsos positivos, falsos negativos e integración de evidencias}

El escaso solapamiento entre los primeros cribados obligó a preguntarse cuánto había de señal y cuánto de ruido. \textcite{c16-vonmering2002} compararon sistemáticamente los grandes conjuntos de datos de levadura entre sí y con un conjunto de referencia de interacciones ya descritas, para medir su exactitud y sus sesgos. La conclusión general, que la interactómica ha confirmado desde entonces, es que cada método tiene su propio perfil de errores: el doble híbrido detecta contactos transitorios pero produce falsos positivos (proteínas ``pegajosas'', autoactivadoras o que nunca coinciden en la célula) y pierde proteínas de membrana; AP-MS captura complejos estables pero confunde asociación indirecta con contacto directo y favorece a las proteínas abundantes. Las interacciones respaldadas por varios métodos independientes son mucho más fiables que las que aparecen en uno solo.

Esa observación sugiere cómo integrar. Si cada fuente de evidencia $c$ asigna a un par una probabilidad $s_c$ de que la asociación sea real, y las fuentes fallan de forma independiente, la probabilidad de que \emph{todas} se equivoquen es el producto de sus errores:
\begin{equation}\label{eq:16-combinada}
S = 1-\prod_{c}\left(1-s_c\right).
\end{equation}
\begin{simbolos}
\simbolo{$s_c$}{Confianza de la fuente o ``canal'' $c$ (experimentos, coexpresión, texto, contexto genómico\dots).}
\simbolo{$S$}{Confianza combinada, bajo el supuesto de independencia entre canales.}
\end{simbolos}
Con tres canales de confianza moderada, $s=(0{,}6;\ 0{,}5;\ 0{,}7)$, la \cref{eq:16-combinada} da $S=1-0{,}4\times0{,}5\times0{,}3=0{,}94$. Esta es la lógica de fondo de \textbf{STRING}, que integra interacciones físicas y asociaciones funcionales procedentes de experimentos, bases de datos curadas, coexpresión, contexto genómico conservado y minería automática de textos, las puntúa y las transfiere a organismos poco estudiados por ortología; su versión 12.0 permite incluso construir la red de un genoma nuevo enviando su proteoma \parencite{c16-szklarczyk2023}. En el otro extremo del espectro, \textbf{BioGRID} no predice nada: cura manualmente interacciones proteicas y genéticas de la literatura primaria, alrededor de 1,93 millones en 2021, además de modificaciones postraduccionales, interacciones con fármacos y cribados CRISPR \parencite{c16-oughtred2021}.

\begin{cuidado}[Una asociación no es un contacto]
Una arista de STRING con puntuación alta puede significar que dos proteínas se tocan, que forman parte del mismo complejo sin tocarse, o simplemente que sus genes se coexpresan o se citan juntos en los artículos. Para preguntas de estructura física, filtre por canal (STRING ofrece una red ``física'') o use datos curados como BioGRID. Y recuerde que la ausencia de arista casi nunca significa ausencia de interacción: los interactomas siguen incompletos y sesgados hacia las proteínas más estudiadas.
\end{cuidado}

\subsection{Tres modelos de red y lo que predicen}

Para saber si una propiedad de una red real es notable, necesitamos compararla con lo que produciría el azar. Tres modelos marcan los hitos de esa comparación.

\paragraph{Erdős-Rényi (ER)} Se unen $n$ nodos colocando cada una de las $\binom{n}{2}$ aristas posibles con probabilidad $p$, de forma independiente. El grado de un nodo es la suma de $n-1$ ensayos de Bernoulli, así que sigue una binomial que, para $n$ grande y $\langle k\rangle=p(n-1)$ fijo, converge a una Poisson:
\begin{equation}\label{eq:16-er}
P(k)=\binom{n-1}{k}p^k(1-p)^{n-1-k}\;\xrightarrow{\;n\to\infty\;}\;e^{-\langle k\rangle}\frac{\langle k\rangle^k}{k!},\qquad C=p,\qquad L\approx\frac{\ln n}{\ln\langle k\rangle}.
\end{equation}
\begin{simbolos}
\simbolo{$p$}{Probabilidad de cada arista.}
\simbolo{$P(k)$}{Distribución de grado: concentrada alrededor de $\langle k\rangle$, con colas que decaen más rápido que exponencialmente.}
\end{simbolos}
El agrupamiento vale $p$ porque dos vecinos de un nodo están unidos, como cualquier otro par, con probabilidad $p$. La fórmula de $L$ se obtiene contando: desde un nodo se alcanzan unos $\langle k\rangle$ nodos a distancia 1, $\langle k\rangle^2$ a distancia 2 y $\langle k\rangle^\ell$ a distancia $\ell$; toda la red queda cubierta cuando $\langle k\rangle^L\approx n$. La distancia crece solo con el \emph{logaritmo} del tamaño.

\paragraph{Watts-Strogatz (WS)} Las redes reales combinan dos rasgos que ER no reproduce a la vez: distancias cortas \emph{y} vecindarios muy agrupados. \textcite{c16-watts1998} mostraron que ambos coexisten en un modelo sencillo: se parte de un anillo en el que cada nodo se une a sus $k$ vecinos más próximos (muy agrupado, pero con $L\approx n/2k$) y se reconecta cada arista al azar con probabilidad $p$. Unos pocos ``atajos'' bastan para desplomar $L$, mientras que $C$ apenas cambia porque destruir triángulos requiere reconectar muchas aristas. Esa franja intermedia es el régimen de \emph{mundo pequeño} (\cref{fig:16-ws}).

\begin{figure}[t]
\centering
\begin{tikzpicture}
\begin{axis}[curso, every axis plot/.append style={thick}, width=0.8\linewidth, height=5.4cm, xmode=log,
   xmin=1e-4, xmax=1, ymin=0, ymax=1.08,
   xlabel={probabilidad de reconexión $p$}, ylabel={valor relativo a $p=0$},
   legend style={at={(0.5,-0.3)},anchor=north}, legend columns=2]
  \fill[amarillo!15] (axis cs:5e-4,0) rectangle (axis cs:3e-2,1.08);
  \addplot[azul, mark=*, mark size=1.6pt] table[x=p,y=C] {figuras/cap16/ws.dat};
  \addlegendentry{$C(p)/C(0)$ agrupamiento}
  \addplot[naranja, mark=square*, mark size=1.5pt] table[x=p,y=L] {figuras/cap16/ws.dat};
  \addlegendentry{$L(p)/L(0)$ distancia media}
  \node[etiqueta, text=amarillo!40!black, align=center] at (axis cs:4e-3,0.55) {mundo\\pequeño};
\end{axis}
\end{tikzpicture}
\caption[El régimen de mundo pequeño]{Modelo de Watts-Strogatz con $n=1000$ nodos y $k=10$ vecinos (promedio de cuatro redes por punto). Con $p=0{,}01$ (uno de cada cien enlaces reconectado) la distancia media ya cayó al 18\,\% de su valor inicial, mientras que el agrupamiento conserva el 97\,\%. La banda amarilla marca el régimen de mundo pequeño: tan agrupado como una red regular y casi tan compacto como una aleatoria. Calculado con NetworkX en \archivo{figuras/cap16/generar.py}.}
\label{fig:16-ws}
\end{figure}

\paragraph{Barabási-Albert (BA)} Ni ER ni WS producen nodos con muchísimas más conexiones que la media. \textcite{c16-barabasi1999} observaron distribuciones de grado de cola larga en varias redes reales y propusieron dos ingredientes para explicarlas: \emph{crecimiento} (la red gana nodos con el tiempo) y \emph{enlace preferencial} (los nodos nuevos prefieren unirse a los ya muy conectados, ``el rico se hace más rico'').

\begin{teorema}{Distribución de grado del modelo BA}{16-ba}
Si en cada paso de tiempo se añade un nodo con $m$ aristas y cada arista se une al nodo $i$ con probabilidad $\Pi(k_i)=k_i/\sum_j k_j$, entonces, en la aproximación de continuo y para $t\to\infty$,
\begin{equation}\label{eq:16-ba}
k_i(t)=m\left(\frac{t}{t_i}\right)^{1/2},\qquad P(k)=\frac{2m^2}{k^{3}}\quad (k\ge m).
\end{equation}
Es decir, una ley de potencias $P(k)\propto k^{-\gamma}$ con $\gamma=3$, independiente de $m$.
\end{teorema}
\begin{simbolos}
\simbolo{$t_i$}{Instante en que el nodo $i$ se incorporó a la red.}
\simbolo{$\Pi(k_i)$}{Probabilidad de que una arista nueva elija al nodo $i$.}
\simbolo{$\gamma$}{Exponente de la ley de potencias.}
\end{simbolos}

La derivación es breve y muy instructiva. En el instante $t$ hay unos $t$ nodos y $mt$ aristas, así que $\sum_j k_j=2mt$. Cada nodo nuevo reparte $m$ aristas, y el nodo $i$ recibe en promedio $m\,\Pi(k_i)$ de ellas:
\[
\frac{dk_i}{dt}=m\,\frac{k_i}{2mt}=\frac{k_i}{2t}
\;\Longrightarrow\; \ln k_i = \tfrac12\ln t + \text{cte}
\;\Longrightarrow\; k_i(t)=m\left(\frac{t}{t_i}\right)^{1/2},
\]
usando la condición inicial $k_i(t_i)=m$. Los nodos más viejos son los más conectados. Como los tiempos de llegada $t_i$ están repartidos uniformemente en $[0,t]$, la probabilidad de que un nodo tenga grado menor que $k$ es
\[
\Prob\big(k_i(t)<k\big)=\Prob\!\left(t_i>\frac{m^2t}{k^2}\right)=1-\frac{m^2}{k^2},
\]
y al derivar respecto a $k$ obtenemos $P(k)=2m^2/k^3$. Una ley de potencias no tiene una escala característica (de ahí el nombre de \emph{red libre de escala}): su varianza diverge cuando $\gamma\le3$, y siempre hay una probabilidad no despreciable de encontrar nodos con grado cientos de veces mayor que la media, los llamados \ingles{hubs}.

\subsection{La red de la levadura frente a los modelos}

¿Qué aspecto tiene una red real? Tomemos la red física de levadura de STRING v12.0 con puntuación combinada de al menos 700 (``confianza alta''). Tiene $n=\num{3384}$ proteínas y $m=\num{43030}$ interacciones, con $\langle k\rangle=25{,}4$ y una componente gigante de \num{2780} nodos. Su coeficiente de agrupamiento es $C=0{,}58$, mientras que una red ER con los mismos $n$ y $m$ tendría $C=p=0{,}0075$: casi ochenta veces menos. La distancia media en la componente gigante es $L\approx5{,}2$ pasos, frente a $2{,}8$ en la red ER equivalente; ambas son minúsculas comparadas con el número de nodos. La red de levadura es, pues, un mundo pequeño muy agrupado.

\begin{figure}[t]
\centering
\begin{tikzpicture}
\begin{axis}[curso, every axis plot/.append style={thick}, width=0.86\linewidth, height=6.4cm, xmode=log, ymode=log,
   xmin=0.8, xmax=500, ymin=1e-7, ymax=1,
   xlabel={grado $k$}, ylabel={$P(k)$ (agrupado logarítmicamente)},
   legend pos=south west]
  \addplot[azul, mark=*, mark size=1.8pt] table[x=k,y=pk] {figuras/cap16/grado_levadura.dat};
  \addlegendentry{levadura, STRING $\ge700$}
  \addplot[aqua, mark=triangle*, mark size=2pt] table[x=k,y=pk] {figuras/cap16/grado_er.dat};
  \addlegendentry{Erdős-Rényi (mismos $n$, $m$)}
  \addplot[naranja, mark=square*, mark size=1.5pt] table[x=k,y=pk] {figuras/cap16/grado_ba.dat};
  \addlegendentry{Barabási-Albert ($m=13$)}
  \addplot[tinta2, dashed, thick, domain=26:400, samples=20] {338/x^3};
  \addlegendentry{teoría BA: $2m^2/k^3$}
\end{axis}
\end{tikzpicture}
\caption[Distribución de grado en escala doble logarítmica]{Distribución de grado de la red física de levadura (STRING v12.0, confianza $\ge700$) comparada con dos modelos del mismo tamaño. La red ER (verde) concentra todos sus nodos alrededor de $\langle k\rangle\approx25$, sin cola. El modelo BA (naranja) sigue la recta de pendiente $-3$ predicha por la \cref{eq:16-ba}. La levadura muestra una cola ancha y extendida, pero no una recta limpia: una meseta hasta $k\approx300$ producida por las proteínas ribosómicas, que forman una ``bola'' casi completa de aristas. Los puntos resultan de agrupar los grados en intervalos de anchura logarítmica creciente y dividir por la anchura de cada uno.}
\label{fig:16-grado}
\end{figure}

La \cref{fig:16-grado} muestra su distribución de grado. A diferencia de la red ER, la levadura tiene una cola larga: proteínas con 200 o más interacciones. Pero la cola no es la recta limpia del modelo BA, y los nodos de mayor grado cuentan una historia menos romántica que ``el rico se hace más rico'': 49 de las 50 proteínas más conectadas son proteínas ribosómicas (RPS13, RPS7A, RPS5\dots). El ribosoma es un complejo enorme purificado muchas veces, y el canal físico de STRING lo representa como un grupo casi completamente conectado. Una estimación por máxima verosimilitud del exponente en la cola ($k\ge26$), $\hat\gamma=1+N\big[\sum_i\ln\big(k_i/(k_{\min}-\tfrac12)\big)\big]^{-1}$, da $1{,}99$ para la levadura y $2{,}83$ para la red BA (cuyo valor teórico es 3). Un exponente tan cercano a 2, el valor por debajo del cual ni siquiera el grado medio estaría definido en una red infinita, delata que la cola está dominada por unos pocos complejos grandes y no por un proceso de crecimiento como el del modelo BA.

\begin{historia}[¿Son raras las redes libres de escala?]
Durante una década, la ``ley de potencias'' fue casi un dogma de la ciencia de redes. \textcite{c16-broido2019} la sometieron a una prueba severa: aplicaron herramientas estadísticas modernas a casi \num{1000} redes sociales, biológicas, tecnológicas, de transporte y de información, y encontraron que las redes \emph{fuertemente} libres de escala son empíricamente raras, que para la mayoría una distribución log-normal se ajusta igual de bien o mejor que una ley de potencias, y que solo un puñado de redes tecnológicas y biológicas parecen fuertemente libres de escala. El debate que siguió fue intenso. La lección práctica es metodológica: una recta aproximada en un gráfico log-log \emph{no} demuestra una ley de potencias; hay que ajustar por máxima verosimilitud, elegir $k_{\min}$ con criterio y comparar con distribuciones alternativas.
\end{historia}

\subsection{Centralidad, esencialidad y medicina de redes}

Si la topología importa, las proteínas centrales deberían ser especiales. \textcite{c16-jeong2001} analizaron la red de interacciones de levadura y encontraron que la probabilidad de que la eliminación de una proteína sea letal para la célula crece con su número de interacciones: las proteínas muy conectadas tienden a ser esenciales. Este principio de \emph{centralidad-letalidad} ha sido matizado después (parte de la correlación se debe a que los complejos esenciales son grandes, justamente el efecto que vimos con el ribosoma), pero la idea de que la posición en la red informa sobre la función se mantiene.

La \emph{medicina de redes} lleva esta idea a la enfermedad humana. \textcite{c16-barabasi2011} argumentan que, dada la interdependencia funcional de los componentes celulares, una enfermedad rara vez es consecuencia de la anomalía de un único gen: refleja la perturbación de una red que conecta tejidos y órganos. Los genes asociados a una misma enfermedad tienden a agruparse en un \emph{módulo de enfermedad}, lo que permite priorizar genes candidatos de un estudio de asociación por su proximidad en la red a genes ya conocidos, o proponer dianas terapéuticas y nuevos usos para fármacos existentes.

\begin{ideaclave}
Grado: cuántos vecinos. Agrupamiento: cuán unidos están esos vecinos. Intermediación: cuánto tráfico pasa por el nodo. Tres números distintos para tres preguntas distintas.
\end{ideaclave}

\subsection{Módulos y comunidades}

Una célula no es una sopa homogénea de interacciones: tiene complejos, vías y orgánulos. En el lenguaje de redes, esto se traduce en \emph{comunidades}: grupos de nodos densamente conectados entre sí y escasamente conectados con el resto. \textcite{c16-girvan2002} propusieron detectarlas eliminando progresivamente las aristas de mayor intermediación, que son las que hacen de puente entre grupos; al retirarlas, la red se fragmenta en sus comunidades naturales. Faltaba un criterio para decidir cuándo detenerse, y ese criterio es la modularidad.

\begin{definicion}{Modularidad}{16-modularidad}
Dada una partición de los nodos en comunidades $c_i$, su modularidad es la fracción de aristas dentro de las comunidades menos la fracción esperada si las aristas se colocaran al azar respetando los grados:
\begin{equation}\label{eq:16-modularidad}
Q=\frac{1}{2m}\sum_{i,j}\left[A_{ij}-\frac{k_ik_j}{2m}\right]\delta(c_i,c_j)
=\sum_{c}\left[\frac{l_c}{m}-\left(\frac{d_c}{2m}\right)^{2}\right].
\end{equation}
\end{definicion}
\begin{simbolos}
\simbolo{$k_ik_j/2m$}{Número esperado de aristas entre $i$ y $j$ en el modelo nulo de configuración (mismos grados, aristas al azar).}
\simbolo{$\delta(c_i,c_j)$}{Vale 1 si $i$ y $j$ están en la misma comunidad.}
\simbolo{$l_c$}{Número de aristas dentro de la comunidad $c$.}
\simbolo{$d_c$}{Suma de los grados de los nodos de $c$.}
\end{simbolos}
El término $k_ik_j/2m$ merece una explicación. Si cortamos cada arista en dos ``medias aristas'' y las volvemos a emparejar al azar, cada una de las $k_i$ medias aristas de $i$ se empareja con alguna de $j$ con probabilidad $k_j/(2m-1)\approx k_j/2m$. Para pasar de la primera forma a la segunda, basta agrupar la suma por comunidades: $\sum_{i,j\in c}A_{ij}=2l_c$ y $\sum_{i,j\in c}k_ik_j=d_c^2$. \textcite{c16-newman2006} mostró que $Q$ se puede escribir con los autovectores de la \emph{matriz de modularidad} $B_{ij}=A_{ij}-k_ik_j/2m$, lo que da un algoritmo espectral eficaz; hoy el método más usado es el de Louvain \parencite{c16-blondel2008}, que mueve nodos entre comunidades mientras $Q$ aumente y luego agrega cada comunidad en un supernodo, repitiendo el proceso hasta converger.

\begin{ejemplo}{Modularidad del grafo de juguete}{16-modularidad}
En la \cref{fig:16-grafo}, con $m=8$, consideremos la partición $\{$A, B, C$\}\,|\,\{$D, E, F, G$\}$. La primera comunidad tiene $l=3$ aristas y grado total $d=2+2+3=7$; la segunda, $l=4$ y $d=3+2+3+1=9$. Por la \cref{eq:16-modularidad},
\[
Q=\left[\tfrac38-\left(\tfrac{7}{16}\right)^2\right]+\left[\tfrac48-\left(\tfrac{9}{16}\right)^2\right]=0{,}1836+0{,}1836=0{,}367.
\]
Si en cambio movemos D al primer grupo, $\{$A, B, C, D$\}\,|\,\{$E, F, G$\}$, la modularidad baja a $0{,}219$: la partición natural corta la única arista puente, C--D.
\end{ejemplo}

La \cref{fig:16-fus3} aplica todo esto a una red real: las proteínas a distancia $\le2$ de la quinasa FUS3 en la red de levadura. El algoritmo de Louvain encuentra cinco comunidades con $Q=0{,}57$ que coinciden con módulos biológicos reconocibles: el núcleo de la cascada MAP quinasa de la feromona (FUS3, STE7, KSS1, el factor de transcripción STE12 y sus inhibidores DIG1 y DIG2), la parte proximal de la vía (la subunidad $\beta\gamma$ de la proteína G, STE4--STE18, el andamio STE5, la quinasa STE20, la GTPasa CDC42), la vía de alta osmolaridad (HOG1, PBS2, SSK2, SSK22), y el control del ciclo celular (CDC28, las ciclinas CLN2 y CLB5, y FAR1, que detiene el ciclo en respuesta a la feromona). La quinasa FUS3 tiene la mayor intermediación porque conecta varios de estos módulos. Para saber si $Q=0{,}57$ es alto, repetimos el análisis en 20 versiones aleatorizadas de la misma red que conservan el grado de cada nodo: dan $Q=0{,}36\pm0{,}02$. La estructura modular es real, no un artefacto del algoritmo.

\begin{figure}[!t]
\centering
\begin{tikzpicture}
\node[circle,draw=white,line width=0.4pt,fill=aqua,minimum size=3.8mm,inner sep=0pt] (n-TIM9) at (10.67,4.79) {};
\node[font=\sffamily\tiny,text=tinta,inner sep=1pt,anchor=191] at (n-TIM9.11) {TIM9};
\node[circle,draw=white,line width=0.4pt,fill=azul,minimum size=3.7mm,inner sep=0pt] (n-STE4) at (2.05,5.60) {};
\node[font=\sffamily\tiny,text=tinta,inner sep=1pt,anchor=334] at (n-STE4.154) {STE4};
\node[circle,draw=white,line width=0.4pt,fill=azul,minimum size=3.5mm,inner sep=0pt] (n-BEM1) at (0.00,4.38) {};
\node[font=\sffamily\tiny,text=tinta,inner sep=1pt,anchor=354] at (n-BEM1.174) {BEM1};
\node[circle,draw=white,line width=0.4pt,fill=magenta,minimum size=3.2mm,inner sep=0pt] (n-CLN2) at (5.23,7.38) {};
\node[font=\sffamily\tiny,text=tinta,inner sep=1pt,anchor=277] at (n-CLN2.97) {CLN2};
\node[circle,draw=white,line width=0.4pt,fill=naranja,minimum size=3.2mm,inner sep=0pt] (n-DIG1) at (10.75,1.98) {};
\node[font=\sffamily\tiny,text=tinta,inner sep=1pt,anchor=160] at (n-DIG1.-20) {DIG1};
\node[circle,draw=white,line width=0.4pt,fill=aqua,minimum size=3.9mm,inner sep=0pt] (n-TIM12) at (8.45,4.99) {};
\node[font=\sffamily\tiny,text=tinta,inner sep=1pt,anchor=203] at (n-TIM12.23) {TIM12};
\node[circle,draw=white,line width=0.4pt,fill=violeta,minimum size=3.0mm,inner sep=0pt] (n-PTP2) at (2.88,0.18) {};
\node[font=\sffamily\tiny,text=tinta,inner sep=1pt,anchor=52] at (n-PTP2.-128) {PTP2};
\node[circle,draw=white,line width=0.4pt,fill=naranja,minimum size=3.0mm,inner sep=0pt] (n-PIN3) at (9.15,1.99) {};
\node[font=\sffamily\tiny,text=tinta,inner sep=1pt,anchor=152] at (n-PIN3.-28) {PIN3};
\node[circle,draw=white,line width=0.4pt,fill=aqua,minimum size=3.8mm,inner sep=0pt] (n-TIM22) at (11.09,6.08) {};
\node[font=\sffamily\tiny,text=tinta,inner sep=1pt,anchor=203] at (n-TIM22.23) {TIM22};
\node[circle,draw=white,line width=0.4pt,fill=azul,minimum size=3.0mm,inner sep=0pt] (n-GDI1) at (1.72,1.44) {};
\node[font=\sffamily\tiny,text=tinta,inner sep=1pt,anchor=31] at (n-GDI1.-149) {GDI1};
\node[circle,draw=white,line width=0.4pt,fill=azul,minimum size=3.5mm,inner sep=0pt] (n-STE18) at (1.84,6.27) {};
\node[font=\sffamily\tiny,text=tinta,inner sep=1pt,anchor=327] at (n-STE18.147) {STE18};
\node[circle,draw=white,line width=0.4pt,fill=azul,minimum size=3.8mm,inner sep=0pt] (n-CDC42) at (1.13,3.36) {};
\node[font=\sffamily\tiny,text=tinta,inner sep=1pt,anchor=5] at (n-CDC42.-175) {CDC42};
\node[circle,draw=white,line width=0.4pt,fill=azul,minimum size=3.5mm,inner sep=0pt] (n-SHO1) at (2.61,1.98) {};
\node[font=\sffamily\tiny,text=tinta,inner sep=1pt,anchor=31] at (n-SHO1.-149) {SHO1};
\node[circle,draw=white,line width=0.4pt,fill=aqua,minimum size=3.8mm,inner sep=0pt] (n-TIM18) at (10.14,6.61) {};
\node[font=\sffamily\tiny,text=tinta,inner sep=1pt,anchor=212] at (n-TIM18.32) {TIM18};
\node[circle,draw=white,line width=0.4pt,fill=naranja,minimum size=3.0mm,inner sep=0pt] (n-MSG5) at (6.43,5.12) {};
\node[font=\sffamily\tiny,text=tinta,inner sep=1pt,anchor=241] at (n-MSG5.61) {MSG5};
\node[circle,draw=white,line width=0.4pt,fill=violeta,minimum size=3.4mm,inner sep=0pt] (n-PTP3) at (4.71,2.07) {};
\node[font=\sffamily\tiny,text=tinta,inner sep=1pt,anchor=60] at (n-PTP3.-120) {PTP3};
\node[circle,draw=white,line width=0.4pt,fill=magenta,minimum size=3.4mm,inner sep=0pt] (n-CDC28) at (3.88,7.60) {};
\node[font=\sffamily\tiny,text=tinta,inner sep=1pt,anchor=296] at (n-CDC28.116) {CDC28};
\node[circle,draw=white,line width=0.4pt,fill=magenta,minimum size=3.2mm,inner sep=0pt] (n-CLB5) at (2.42,7.42) {};
\node[font=\sffamily\tiny,text=tinta,inner sep=1pt,anchor=312] at (n-CLB5.132) {CLB5};
\node[circle,draw=white,line width=0.4pt,fill=aqua,minimum size=3.8mm,inner sep=0pt] (n-SDH3) at (9.14,6.39) {};
\node[font=\sffamily\tiny,text=tinta,inner sep=1pt,anchor=217] at (n-SDH3.37) {SDH3};
\node[circle,draw=white,line width=0.4pt,fill=magenta,minimum size=3.9mm,inner sep=0pt] (n-FAR1) at (4.04,5.78) {};
\node[font=\sffamily\tiny,text=tinta,inner sep=1pt,anchor=310] at (n-FAR1.130) {FAR1};
\node[circle,draw=white,line width=0.4pt,fill=azul,minimum size=3.8mm,inner sep=0pt] (n-STE5) at (3.37,4.51) {};
\node[font=\sffamily\tiny,text=tinta,inner sep=1pt,anchor=343] at (n-STE5.163) {STE5};
\node[circle,draw=white,line width=0.4pt,fill=violeta,minimum size=3.2mm,inner sep=0pt] (n-SSK2) at (7.20,0.44) {};
\node[font=\sffamily\tiny,text=tinta,inner sep=1pt,anchor=114] at (n-SSK2.-66) {SSK2};
\node[circle,draw=white,line width=0.4pt,fill=azul,minimum size=3.4mm,inner sep=0pt] (n-STE50) at (0.90,2.32) {};
\node[font=\sffamily\tiny,text=tinta,inner sep=1pt,anchor=17] at (n-STE50.-163) {STE50};
\node[circle,draw=white,line width=0.4pt,fill=violeta,minimum size=3.2mm,inner sep=0pt] (n-HOG1) at (4.54,0.39) {};
\node[font=\sffamily\tiny,text=tinta,inner sep=1pt,anchor=71] at (n-HOG1.-109) {HOG1};
\node[circle,draw=white,line width=0.4pt,fill=naranja,minimum size=3.5mm,inner sep=0pt] (n-STE7) at (4.78,3.80) {};
\node[font=\sffamily\tiny,text=tinta,inner sep=1pt,anchor=360] at (n-STE7.180) {STE7};
\node[circle,draw=white,line width=0.4pt,fill=naranja,minimum size=3.4mm,inner sep=0pt] (n-STE12) at (8.52,1.44) {};
\node[font=\sffamily\tiny,text=tinta,inner sep=1pt,anchor=140] at (n-STE12.-40) {STE12};
\node[circle,draw=white,line width=0.4pt,fill=naranja,minimum size=3.7mm,inner sep=0pt] (n-KSS1) at (6.51,2.82) {};
\node[font=\sffamily\tiny,text=tinta,inner sep=1pt,anchor=130] at (n-KSS1.-50) {KSS1};
\node[circle,draw=white,line width=0.4pt,fill=aqua,minimum size=3.8mm,inner sep=0pt] (n-TIM54) at (9.93,5.60) {};
\node[font=\sffamily\tiny,text=tinta,inner sep=1pt,anchor=203] at (n-TIM54.23) {TIM54};
\node[circle,draw=white,line width=0.4pt,fill=aqua,minimum size=3.8mm,inner sep=0pt] (n-TIM10) at (11.40,5.29) {};
\node[font=\sffamily\tiny,text=tinta,inner sep=1pt,anchor=195] at (n-TIM10.15) {TIM10};
\node[circle,draw=white,line width=0.4pt,fill=azul,minimum size=3.8mm,inner sep=0pt] (n-CDC24) at (0.80,5.27) {};
\node[font=\sffamily\tiny,text=tinta,inner sep=1pt,anchor=343] at (n-CDC24.163) {CDC24};
\node[circle,draw=white,line width=0.4pt,fill=azul,minimum size=3.7mm,inner sep=0pt] (n-STE20) at (1.66,4.12) {};
\node[font=\sffamily\tiny,text=tinta,inner sep=1pt,anchor=355] at (n-STE20.175) {STE20};
\node[circle,draw=white,line width=0.4pt,fill=violeta,minimum size=3.8mm,inner sep=0pt] (n-PBS2) at (5.42,1.96) {};
\node[font=\sffamily\tiny,text=tinta,inner sep=1pt,anchor=81] at (n-PBS2.-99) {PBS2};
\node[circle,draw=white,line width=0.4pt,fill=naranja,minimum size=3.4mm,inner sep=0pt] (n-DIG2) at (8.26,3.07) {};
\node[font=\sffamily\tiny,text=tinta,inner sep=1pt,anchor=164] at (n-DIG2.-16) {DIG2};
\node[circle,draw=white,line width=0.4pt,fill=azul,minimum size=4.2mm,inner sep=0pt] (n-STE11) at (3.67,3.04) {};
\node[font=\sffamily\tiny,text=tinta,inner sep=1pt,anchor=21] at (n-STE11.-159) {STE11};
\node[circle,draw=white,line width=0.4pt,fill=violeta,minimum size=3.2mm,inner sep=0pt] (n-SSK22) at (5.79,0.00) {};
\node[font=\sffamily\tiny,text=tinta,inner sep=1pt,anchor=91] at (n-SSK22.-89) {SSK22};
\node[circle,draw=white,line width=0.4pt,fill=naranja,minimum size=4.2mm,inner sep=0pt] (n-FUS3) at (5.88,3.94) {};
\node[font=\sffamily\tiny,text=tinta,inner sep=1pt,anchor=218] at (n-FUS3.38) {FUS3};
\begin{scope}[on background layer]
\draw[rejilla!70!gris,line width=0.35pt] (n-TIM9) -- (n-TIM12);
\draw[rejilla!70!gris,line width=0.35pt] (n-TIM9) -- (n-TIM22);
\draw[rejilla!70!gris,line width=0.35pt] (n-TIM9) -- (n-TIM54);
\draw[rejilla!70!gris,line width=0.35pt] (n-TIM9) -- (n-SDH3);
\draw[rejilla!70!gris,line width=0.35pt] (n-TIM9) -- (n-TIM10);
\draw[rejilla!70!gris,line width=0.35pt] (n-TIM9) -- (n-TIM18);
\draw[rejilla!70!gris,line width=0.35pt] (n-STE4) -- (n-CDC24);
\draw[rejilla!70!gris,line width=0.35pt] (n-STE4) -- (n-STE5);
\draw[rejilla!70!gris,line width=0.35pt] (n-STE4) -- (n-STE20);
\draw[tinta2,line width=0.7pt,densely dashed] (n-STE4) -- (n-FAR1);
\draw[rejilla!70!gris,line width=0.35pt] (n-STE4) -- (n-STE18);
\draw[rejilla!70!gris,line width=0.35pt] (n-BEM1) -- (n-CDC24);
\draw[rejilla!70!gris,line width=0.35pt] (n-BEM1) -- (n-STE5);
\draw[rejilla!70!gris,line width=0.35pt] (n-BEM1) -- (n-CDC42);
\draw[rejilla!70!gris,line width=0.35pt] (n-BEM1) -- (n-STE20);
\draw[rejilla!70!gris,line width=0.35pt] (n-CLN2) -- (n-CDC28);
\draw[rejilla!70!gris,line width=0.35pt] (n-CLN2) -- (n-FAR1);
\draw[rejilla!70!gris,line width=0.35pt] (n-DIG1) -- (n-DIG2);
\draw[rejilla!70!gris,line width=0.35pt] (n-DIG1) -- (n-STE12);
\draw[tinta2,line width=0.7pt,densely dashed] (n-TIM12) -- (n-FUS3);
\draw[rejilla!70!gris,line width=0.35pt] (n-TIM12) -- (n-TIM22);
\draw[rejilla!70!gris,line width=0.35pt] (n-TIM12) -- (n-TIM10);
\draw[rejilla!70!gris,line width=0.35pt] (n-TIM12) -- (n-SDH3);
\draw[rejilla!70!gris,line width=0.35pt] (n-TIM12) -- (n-TIM54);
\draw[rejilla!70!gris,line width=0.35pt] (n-TIM12) -- (n-TIM18);
\draw[rejilla!70!gris,line width=0.35pt] (n-PTP2) -- (n-PTP3);
\draw[rejilla!70!gris,line width=0.35pt] (n-PIN3) -- (n-KSS1);
\draw[rejilla!70!gris,line width=0.35pt] (n-TIM22) -- (n-SDH3);
\draw[rejilla!70!gris,line width=0.35pt] (n-TIM22) -- (n-TIM10);
\draw[rejilla!70!gris,line width=0.35pt] (n-TIM22) -- (n-TIM54);
\draw[rejilla!70!gris,line width=0.35pt] (n-TIM22) -- (n-TIM18);
\draw[rejilla!70!gris,line width=0.35pt] (n-GDI1) -- (n-STE11);
\draw[rejilla!70!gris,line width=0.35pt] (n-STE18) -- (n-CDC24);
\draw[rejilla!70!gris,line width=0.35pt] (n-STE18) -- (n-STE5);
\draw[tinta2,line width=0.7pt,densely dashed] (n-STE18) -- (n-FAR1);
\draw[rejilla!70!gris,line width=0.35pt] (n-CDC42) -- (n-CDC24);
\draw[rejilla!70!gris,line width=0.35pt] (n-CDC42) -- (n-STE50);
\draw[rejilla!70!gris,line width=0.35pt] (n-CDC42) -- (n-SHO1);
\draw[rejilla!70!gris,line width=0.35pt] (n-CDC42) -- (n-STE20);
\draw[rejilla!70!gris,line width=0.35pt] (n-CDC42) -- (n-STE11);
\draw[rejilla!70!gris,line width=0.35pt] (n-SHO1) -- (n-STE50);
\draw[tinta2,line width=0.7pt,densely dashed] (n-SHO1) -- (n-PBS2);
\draw[rejilla!70!gris,line width=0.35pt] (n-SHO1) -- (n-STE11);
\draw[rejilla!70!gris,line width=0.35pt] (n-TIM18) -- (n-TIM10);
\draw[rejilla!70!gris,line width=0.35pt] (n-TIM18) -- (n-TIM54);
\draw[rejilla!70!gris,line width=0.35pt] (n-TIM18) -- (n-SDH3);
\draw[rejilla!70!gris,line width=0.35pt] (n-MSG5) -- (n-FUS3);
\draw[tinta2,line width=0.7pt,densely dashed] (n-PTP3) -- (n-FUS3);
\draw[rejilla!70!gris,line width=0.35pt] (n-PTP3) -- (n-HOG1);
\draw[rejilla!70!gris,line width=0.35pt] (n-CDC28) -- (n-FAR1);
\draw[rejilla!70!gris,line width=0.35pt] (n-CDC28) -- (n-CLB5);
\draw[rejilla!70!gris,line width=0.35pt] (n-CLB5) -- (n-FAR1);
\draw[rejilla!70!gris,line width=0.35pt] (n-SDH3) -- (n-TIM10);
\draw[rejilla!70!gris,line width=0.35pt] (n-SDH3) -- (n-TIM54);
\draw[tinta2,line width=0.7pt,densely dashed] (n-FAR1) -- (n-CDC24);
\draw[tinta2,line width=0.7pt,densely dashed] (n-FAR1) -- (n-FUS3);
\draw[tinta2,line width=0.7pt,densely dashed] (n-STE5) -- (n-FUS3);
\draw[tinta2,line width=0.7pt,densely dashed] (n-STE5) -- (n-STE7);
\draw[rejilla!70!gris,line width=0.35pt] (n-STE5) -- (n-STE11);
\draw[rejilla!70!gris,line width=0.35pt] (n-SSK2) -- (n-SSK22);
\draw[rejilla!70!gris,line width=0.35pt] (n-SSK2) -- (n-PBS2);
\draw[rejilla!70!gris,line width=0.35pt] (n-STE50) -- (n-STE11);
\draw[rejilla!70!gris,line width=0.35pt] (n-HOG1) -- (n-PBS2);
\draw[rejilla!70!gris,line width=0.35pt] (n-STE7) -- (n-FUS3);
\draw[tinta2,line width=0.7pt,densely dashed] (n-STE7) -- (n-STE11);
\draw[rejilla!70!gris,line width=0.35pt] (n-STE7) -- (n-KSS1);
\draw[rejilla!70!gris,line width=0.35pt] (n-STE12) -- (n-DIG2);
\draw[rejilla!70!gris,line width=0.35pt] (n-STE12) -- (n-KSS1);
\draw[rejilla!70!gris,line width=0.35pt] (n-KSS1) -- (n-FUS3);
\draw[tinta2,line width=0.7pt,densely dashed] (n-KSS1) -- (n-STE11);
\draw[rejilla!70!gris,line width=0.35pt] (n-TIM54) -- (n-TIM10);
\draw[rejilla!70!gris,line width=0.35pt] (n-CDC24) -- (n-STE20);
\draw[rejilla!70!gris,line width=0.35pt] (n-STE20) -- (n-STE11);
\draw[tinta2,line width=0.7pt,densely dashed] (n-PBS2) -- (n-FUS3);
\draw[rejilla!70!gris,line width=0.35pt] (n-PBS2) -- (n-SSK22);
\draw[tinta2,line width=0.7pt,densely dashed] (n-PBS2) -- (n-STE11);
\draw[rejilla!70!gris,line width=0.35pt] (n-DIG2) -- (n-FUS3);
\draw[tinta2,line width=0.7pt,densely dashed] (n-STE11) -- (n-FUS3);
\end{scope}

\begin{scope}[shift={(0.3,-0.9)},font=\sffamily\scriptsize]
  \foreach \c/\t/\x/\y in {azul/{vía proximal: G$\beta\gamma$, STE5, CDC42}/0/0,
      naranja/{cascada MAPK: FUS3, STE7, STE12}/5.6/0,
      aqua/{importación mitocondrial (TIM)}/0/-0.42,
      violeta/{vía HOG (osmolaridad)}/5.6/-0.42,
      magenta/{ciclo celular: CDC28, FAR1}/0/-0.84}{
    \fill[\c] (\x,\y) circle (1.3mm); \node[anchor=west,text=tinta2] at (\x+0.15,\y) {\t};}
  \draw[tinta2,line width=0.7pt,densely dashed] (5.47,-0.84) -- (6.0,-0.84) node[right,text=tinta2]{arista entre comunidades};
\end{scope}
\end{tikzpicture}
\caption[Comunidades en la vecindad de FUS3]{Subred de la vía de la feromona en levadura: todas las proteínas a distancia $\le2$ de FUS3 en la red física de STRING v12.0 con confianza $\ge700$ (36 proteínas, 78 interacciones). Las coordenadas se calcularon con el algoritmo de Kamada-Kawai y los colores son las cinco comunidades de Louvain ($Q=0{,}57$; en redes aleatorizadas con los mismos grados, $Q=0{,}36\pm0{,}02$). El tamaño de cada nodo crece con su grado; las líneas discontinuas son aristas entre comunidades. FUS3 es el nodo de mayor intermediación. Observe el módulo verde de proteínas mitocondriales TIM: llega a la figura por una única arista, FUS3--TIM12, un recordatorio de que una red extraída por vecindad hereda cualquier interacción dudosa de su nodo semilla. Generado con \archivo{figuras/cap16/generar.py}.}
\label{fig:16-fus3}
\end{figure}

\begin{cuidado}[Comunidades a medida del observador]
La maximización de la modularidad tiene limitaciones conocidas: las particiones con $Q$ casi máximo pueden ser muy distintas entre sí, el algoritmo de Louvain depende del orden aleatorio en que visita los nodos (fije siempre la semilla) y existe un límite de resolución por debajo del cual las comunidades pequeñas se funden con sus vecinas. Además, una ego-red como la de la \cref{fig:16-fus3} está sesgada por construcción hacia su semilla. Trate las comunidades como hipótesis que hay que contrastar con anotaciones funcionales (Gene Ontology, complejos conocidos), no como descubrimientos.
\end{cuidado}

\subsection{Herramientas: NetworkX y Cytoscape}

Para el análisis programático, la biblioteca de referencia en Python es NetworkX \parencite{c16-hagberg2008}, que ofrece estructuras de grafo flexibles y cientos de algoritmos (caminos, centralidades, comunidades, modelos aleatorios). Para la exploración visual e interactiva, y para integrar la red con datos de expresión o anotaciones, el estándar es Cytoscape \parencite{c16-shannon2003}, un entorno de código abierto ampliable con extensiones, entre ellas una que consulta STRING directamente. El siguiente fragmento reproduce los números de esta sección.

\begin{codigo}[red\_levadura.py]
import networkx as nx
import pandas as pd

aristas = pd.read_csv("levadura_string700.tsv.gz", sep="\t",
                      comment="#")
G = nx.from_pandas_edgelist(aristas, "prot1", "prot2", "score")
print(G.number_of_nodes(), G.number_of_edges())   # 3384 43030
print(nx.average_clustering(G))                   # 0.581

E = nx.ego_graph(G, "FUS3", radius=2)             # vecindad de FUS3
com = nx.community.louvain_communities(E, seed=1)
print(len(com), nx.community.modularity(E, com))
bc = nx.betweenness_centrality(E)
print(max(bc, key=bc.get))                        # FUS3
\end{codigo}

% ---------------------------------------------------------------------
\section{Redes de regulación génica y modelos dinámicos}\label{sec:16-grn}
% ---------------------------------------------------------------------
\enlacerepo{16.2 · Regulación génica y modelos ODE}

\begin{intuicion}[El termostato y la bañera]
Piense en una bañera con el grifo abierto y el desagüe destapado. El nivel del agua sube hasta que lo que entra por el grifo iguala lo que sale por el desagüe, y ahí se estabiliza. Si queremos llenarla más rápido sin cambiar el nivel final, hay un truco: abrir el grifo a tope al principio y cerrarlo a medida que el agua se acerca a la marca, como hace un termostato. La concentración de una proteína en una célula se comporta exactamente como esa bañera: se sintetiza (el grifo), se degrada y se diluye al dividirse la célula (el desagüe), y un gen que reprime su propio promotor es un termostato molecular. Con dos o tres ``bañeras'' que se abren y cierran los grifos unas a otras, la célula construye memorias, relojes y filtros.
\end{intuicion}

Las redes de la sección anterior eran mapas estáticos. Las redes de regulación génica, en cambio, son \emph{dirigidas}: sus nodos son genes, y una arista $X\to Y$ significa que el producto del gen $X$, un factor de transcripción, modifica la tasa de transcripción de $Y$ (activándola, $X\to Y$, o reprimiéndola, $X\dashv Y$). Para saber qué hace un circuito así no basta con su topología: necesitamos ecuaciones que describan cómo cambian las concentraciones en el tiempo. El libro de \textcite{c16-alon2019} desarrolla este programa con detalle y es la referencia natural para lo que sigue.

\subsection{De la unión al promotor a la función de Hill}

La cinética enzimática ofrece el primer modelo de respuesta saturable. En la reacción $E+S\rightleftharpoons ES\to E+P$, si el complejo $ES$ alcanza rápidamente un estado cuasiestacionario, la velocidad de formación de producto es
\begin{equation}\label{eq:16-mm}
v=\frac{V_{\max}\,[S]}{K_M+[S]},
\end{equation}
\begin{simbolos}
\simbolo{$v$}{Velocidad de la reacción.}
\simbolo{$V_{\max}$}{Velocidad máxima, cuando la enzima está saturada.}
\simbolo{$K_M$}{Constante de Michaelis: concentración de sustrato que da la mitad de $V_{\max}$.}
\end{simbolos}
la ley que Michaelis y Menten publicaron en 1913 y cuya traducción comentada al inglés publicaron \textcite{c16-johnson2011}. La misma forma aparece cuando un factor de transcripción $X$ se une a un sitio del promotor $D$: en equilibrio, $[D][X]/[DX]=K$, y la fracción de promotores ocupados es $[X]/(K+[X])$.

Muchos factores de transcripción actúan como dímeros o tetrámeros, o se unen a varios sitios de forma cooperativa. Si el promotor solo se activa cuando se unen $n$ moléculas a la vez ($D+nX\rightleftharpoons DX_n$, con $K^n=[D][X]^n/[DX_n]$), la fracción ocupada es
\[
\theta=\frac{[DX_n]}{[D]+[DX_n]}=\frac{[X]^n/K^n}{1+[X]^n/K^n}.
\]

\begin{definicion}{Funciones de Hill}{16-hill}
La tasa de producción de un gen regulado por $X$ se modela como
\begin{equation}\label{eq:16-hill}
f_{\text{act}}(X)=\beta\,\frac{X^n}{K^n+X^n},\qquad
f_{\text{rep}}(X)=\beta\,\frac{1}{1+(X/K)^n}.
\end{equation}
\end{definicion}
\begin{simbolos}
\simbolo{$\beta$}{Tasa máxima de producción (fuerza del promotor).}
\simbolo{$K$}{Coeficiente de activación o represión: concentración de $X$ que da la mitad del efecto.}
\simbolo{$n$}{Coeficiente de Hill: mide la cooperatividad y, con ella, lo abrupta que es la respuesta.}
\end{simbolos}

¿Cuán abrupta? Llamemos $x=X/K$. Para pasar del 10\,\% al 90\,\% de activación hay que resolver $x^n/(1+x^n)=0{,}1$ y $0{,}9$, es decir, $x_{10}=9^{-1/n}$ y $x_{90}=9^{1/n}$. Su cociente,
\begin{equation}\label{eq:16-hill81}
\frac{X_{90}}{X_{10}}=81^{1/n},
\end{equation}
vale 81 para $n=1$, 9 para $n=2$ y solo 3 para $n=4$. Con $n$ grande la función de Hill se aproxima a un escalón, y el gen se comporta como un interruptor lógico: encendido por encima de $K$, apagado por debajo (\cref{fig:16-hill}). Esta es la base de las aproximaciones ``lógicas'' que usaremos para razonar sobre circuitos.

\begin{figure}[t]
\centering
\begin{tikzpicture}
\begin{groupplot}[group style={group size=2 by 1, horizontal sep=1.5cm},
   curso, every axis plot/.append style={thick}, width=0.52\linewidth, height=5.2cm, xmin=0, xmax=4, ymin=0, ymax=1.05,
   xlabel={$X/K$}, domain=0:4, samples=160]
\nextgroupplot[title={activador}, ylabel={producción $f/\beta$}, legend pos=south east]
  \addplot[azul] {x/(1+x)}; \addlegendentry{$n=1$}
  \addplot[naranja] {x^2/(1+x^2)}; \addlegendentry{$n=2$}
  \addplot[violeta] {x^4/(1+x^4)}; \addlegendentry{$n=4$}
  \addplot[tinta2, densely dotted, thick] coordinates {(1,0) (1,0.5) (0,0.5)};
\nextgroupplot[title={represor}, legend pos=north east]
  \addplot[azul] {1/(1+x)}; \addlegendentry{$n=1$}
  \addplot[naranja] {1/(1+x^2)}; \addlegendentry{$n=2$}
  \addplot[violeta] {1/(1+x^4)}; \addlegendentry{$n=4$}
  \addplot[tinta2, densely dotted, thick] coordinates {(1,0) (1,0.5) (0,0.5)};
\end{groupplot}
\end{tikzpicture}
\caption[Funciones de Hill]{Funciones de Hill para un activador (izquierda) y un represor (derecha), \cref{eq:16-hill}. Todas pasan por la mitad del efecto en $X=K$ (líneas punteadas), pero la cooperatividad cambia la forma: con $n=1$ la respuesta es gradual y hace falta multiplicar $X$ por 81 para pasar del 10\,\% al 90\,\%; con $n=4$ basta un factor 3 (\cref{eq:16-hill81}). La pendiente en $X=K$, en escala logarítmica, es $n/4$.}
\label{fig:16-hill}
\end{figure}

\subsection{El modelo básico: producción y eliminación}

Consideremos un gen $Y$ cuyo activador $X$ pasa de golpe, en $t=0$, a un nivel que satura su promotor. A partir de ese momento $Y$ se produce a tasa constante $\beta$ y se elimina en proporción a su concentración:
\begin{equation}\label{eq:16-ode}
\frac{dY}{dt}=\beta-\alpha Y,\qquad \alpha=\alpha_{\text{deg}}+\alpha_{\text{dil}}.
\end{equation}
\begin{simbolos}
\simbolo{$Y(t)$}{Concentración de la proteína $Y$.}
\simbolo{$\alpha_{\text{deg}}$}{Tasa de degradación activa (proteólisis).}
\simbolo{$\alpha_{\text{dil}}$}{Tasa de dilución por crecimiento: $\ln2/\tau$ si la célula se divide cada $\tau$.}
\end{simbolos}

\begin{teorema}{Respuesta de la regulación simple}{16-respuesta}
La solución de la \cref{eq:16-ode} con $Y(0)=0$ es
\begin{equation}\label{eq:16-sol}
Y(t)=Y_{\text{est}}\left(1-e^{-\alpha t}\right),\qquad Y_{\text{est}}=\frac{\beta}{\alpha},\qquad t_{1/2}=\frac{\ln 2}{\alpha}.
\end{equation}
El tiempo de respuesta $t_{1/2}$, en el que $Y$ alcanza la mitad de su valor final, depende solo de la tasa de eliminación, no de la de producción. Al apagar el gen, $Y(t)=Y_{\text{est}}e^{-\alpha t}$ decae con el mismo $t_{1/2}$.
\end{teorema}
La demostración es directa: con $u=Y_{\text{est}}-Y$, la ecuación queda $du/dt=-\alpha u$, cuya solución es $u(t)=u(0)e^{-\alpha t}=Y_{\text{est}}e^{-\alpha t}$. Imponer $Y(t_{1/2})=Y_{\text{est}}/2$ da $e^{-\alpha t_{1/2}}=1/2$.

La consecuencia es contraintuitiva y profunda. Una célula que quiere producir \emph{más} proteína puede usar un promotor más fuerte; pero si quiere producirla \emph{más rápido}, un promotor fuerte no sirve: $t_{1/2}$ no depende de $\beta$. Para las proteínas estables de bacterias, que solo se eliminan por dilución, $\alpha=\ln2/\tau$ y el tiempo de respuesta es exactamente un tiempo de generación.

\begin{ejemplo}{Cuánto tarda en responder un gen}{16-respuesta}
En \especie{Escherichia coli} creciendo con $\tau=30$ min, una proteína estable tiene $\alpha=\ln2/30=0{,}0231\ \text{min}^{-1}$ y $t_{1/2}=30$ min. Si la proteína además se degrada activamente con una vida media de 5 min, $\alpha=0{,}0231+\ln2/5=0{,}162\ \text{min}^{-1}$ y $t_{1/2}=4{,}3$ min, siete veces más rápido; con una vida media de 60 min, $t_{1/2}=20$ min. La degradación acelera la respuesta, pero a un costo: para mantener el mismo nivel final hay que producir la proteína (y gastar energía) a una tasa proporcionalmente mayor.
\end{ejemplo}

\subsection{Autorregulación negativa: acelerar sin desperdiciar}

Existe una forma más económica de acelerar. Supongamos que $X$ reprime su propio promotor:
\begin{equation}\label{eq:16-nar}
\frac{dX}{dt}=\frac{\beta}{1+(X/K)^n}-\alpha X.
\end{equation}
Al principio, con $X\ll K$, el promotor funciona a pleno rendimiento y $X$ crece a la tasa máxima $\beta$. Cuando $X$ se acerca a $K$, la represión frena la producción, y el sistema se estabiliza. Si elegimos un promotor muy fuerte ($\beta$ grande) y una represión que lo compense, obtenemos el mismo nivel final que con regulación simple pero con un arranque mucho más rápido: es el termostato de la bañera. Además, en el límite de represión fuerte ($X_{\text{est}}\gg K$) el estado estacionario satisface $\beta(K/X)^n\approx\alpha X$, es decir,
\begin{equation}\label{eq:16-narest}
X_{\text{est}}\approx\left(\frac{\beta K^n}{\alpha}\right)^{1/(n+1)},
\end{equation}
\begin{simbolos}
\simbolo{$X_{\text{est}}$}{Nivel estacionario de la proteína autorreprimida.}
\simbolo{$n$}{Coeficiente de Hill de la autorrepresión.}
\end{simbolos}
de modo que el nivel depende de $\beta$ solo a través de una raíz de orden $n+1$: si una fluctuación duplica la fuerza del promotor, con $n=2$ el nivel sube un $2^{1/3}-1=26\,\%$, no un $100\,\%$. La autorregulación negativa acelera la respuesta \emph{y} amortigua el ruido en la producción.

\textcite{c16-rosenfeld2002} midieron el efecto con circuitos sintéticos en \especie{E. coli}: la unidad sin autorregulación tarda un ciclo celular en alcanzar la mitad de su nivel final, mientras que la autorregulada lo hace en aproximadamente una quinta parte de ciclo, en acuerdo con la solución analítica del modelo. También señalaron que este motivo aparece en más del 40\,\% de los factores de transcripción conocidos de \especie{E. coli}.

\begin{figure}[t]
\centering
\begin{tikzpicture}
\begin{axis}[curso, every axis plot/.append style={thick}, width=0.84\linewidth, height=5.4cm, xmin=0, xmax=4, ymin=0, ymax=1.08,
   xlabel={tiempo $\alpha t$ (en unidades de $1/\alpha$)}, ylabel={$X/X_{\text{est}}$}, legend pos=south east]
  \addplot[azul] table[x=t,y=simple] {figuras/cap16/nar.dat};
  \addlegendentry{regulación simple, $\beta$}
  \addplot[naranja] table[x=t,y=nar] {figuras/cap16/nar.dat};
  \addlegendentry{autorregulación negativa, $\beta'=12{,}1\beta$}
  \addplot[tinta2, densely dotted, thick] coordinates {(0,0.5) (0.693,0.5) (0.693,0)};
  \addplot[tinta2, densely dotted, thick] coordinates {(0.0855,0) (0.0855,0.5)};
  \node[etiqueta, anchor=west] at (axis cs:0.72,0.12) {$t_{1/2}=\ln 2/\alpha$};
  \node[etiqueta, anchor=west, text=naranja!70!black] at (axis cs:0.11,0.93) {$t_{1/2}=0{,}086/\alpha$};
\end{axis}
\end{tikzpicture}
\caption[La autorregulación negativa acelera la respuesta]{Respuesta a la activación en $t=0$ de un gen con regulación simple (\cref{eq:16-sol}) y de uno que reprime su propio promotor (\cref{eq:16-nar}, con $K=0{,}3X_{\text{est}}$ y $n=2$). Ambos llegan al mismo nivel final, pero el autorregulado usa un promotor 12,1 veces más fuerte que queda frenado a medida que $X$ se acerca a $K$. El resultado es un tiempo de respuesta 8,1 veces menor sin cambiar el estado estacionario ni la tasa de eliminación.}
\label{fig:16-nar}
\end{figure}

\subsection{Motivos de red}

¿Hay patrones que se repiten en las redes de regulación? \textcite{c16-milo2002} definieron los \emph{motivos de red} como patrones de interconexión que aparecen significativamente más a menudo que en redes aleatorizadas, y los encontraron en redes de bioquímica, neurobiología, ecología e ingeniería. Para cada subgrafo de tres o cuatro nodos se cuenta su número de apariciones $N_{\text{real}}$ y se compara con el de un conjunto de redes aleatorizadas que conservan el grado de entrada y de salida de cada nodo:
\begin{equation}\label{eq:16-zscore}
Z=\frac{N_{\text{real}}-\langle N_{\text{azar}}\rangle}{\sigma_{\text{azar}}}.
\end{equation}
\begin{simbolos}
\simbolo{$N_{\text{real}}$}{Apariciones del subgrafo en la red real.}
\simbolo{$\langle N_{\text{azar}}\rangle,\ \sigma_{\text{azar}}$}{Media y desviación estándar en las redes aleatorizadas.}
\simbolo{$Z$}{Puntuación de significación del motivo.}
\end{simbolos}
\textcite{c16-shenorr2002} aplicaron la idea a la red de interacciones transcripcionales directas de \especie{E. coli} y encontraron que buena parte de ella está compuesta por apariciones repetidas de tres motivos muy significativos: el \ingles{feed-forward loop} (FFL), el módulo de entrada única (\ingles{single-input module}, SIM) y los regulones densos superpuestos. Cada motivo tiene una función específica en la dinámica de la expresión. \textcite{c16-alon2007} revisa la teoría y los experimentos que las han puesto a prueba, desde bacterias hasta humanos (\cref{fig:16-motivos}).

\begin{figure}[t]
\centering
\begin{tikzpicture}[font=\sffamily\small,
   g/.style={circle,draw=white,thick,fill=#1,minimum size=7mm,inner sep=0pt,text=white,font=\sffamily\bfseries\small},
   act/.style={-{Stealth[length=2.4mm]},thick,tinta2},
   rep/.style={-{Bar[width=2.6mm]},thick,rojooscuro},
   titulo/.style={font=\sffamily\bfseries\scriptsize,text=tinta,align=center}]
  % NAR
  \node[g=azul] (x1) at (0,0.9) {X};
  \path (x1) edge[rep, out=55, in=125, looseness=7] (x1);
  \node[titulo] at (0.2,-0.5) {autorregulación\\negativa};
  % FFL
  \begin{scope}[xshift=2.6cm]
  \node[g=azul] (fx) at (0,1.6) {X};
  \node[g=naranja] (fy) at (0.9,0.8) {Y};
  \node[g=aqua] (fz) at (0,0) {Z};
  \draw[act] (fx) -- (fy); \draw[act] (fy) -- (fz); \draw[act] (fx) -- (fz);
  \node[etiqueta,text=tinta2] at (-0.55,0.8) {AND};
  \node[titulo] at (0.4,-0.7) {FFL coherente\\tipo 1};
  \end{scope}
  % SIM
  \begin{scope}[xshift=5.0cm]
  \node[g=azul] (sx) at (0.9,1.6) {X};
  \foreach \i/\x in {1/0,2/0.9,3/1.8}{\node[g=aqua,minimum size=6mm] (sz\i) at (\x,0) {Z\textsubscript{\i}}; \draw[act] (sx) -- (sz\i);}
  \node[titulo] at (0.9,-0.7) {entrada única\\(SIM)};
  \end{scope}
  % toggle
  \begin{scope}[xshift=8.1cm]
  \node[g=azul] (tu) at (0,0.8) {U};
  \node[g=naranja] (tv) at (1.5,0.8) {V};
  \draw[rep] (tu) to[bend left=35] (tv);
  \draw[rep] (tv) to[bend left=35] (tu);
  \node[titulo] at (0.75,-0.7) {interruptor\\biestable};
  \end{scope}
  % repressilator
  \begin{scope}[xshift=11.2cm,yshift=0.75cm]
  \node[g=azul] (r1) at (90:0.85) {\scriptsize LacI};
  \node[g=naranja] (r2) at (-30:0.85) {\scriptsize TetR};
  \node[g=aqua] (r3) at (210:0.85) {\scriptsize $\lambda$cI};
  \draw[rep] (r1) -- (r2); \draw[rep] (r2) -- (r3); \draw[rep] (r3) -- (r1);
  \node[titulo] at (0,-1.45) {\ingles{repressilator}\\ };
  \end{scope}
\end{tikzpicture}
\caption[Motivos y circuitos de regulación]{Los circuitos que estudiaremos. Las flechas indican activación y las barras rojas, represión. De izquierda a derecha: la autorregulación negativa acelera la respuesta (\cref{fig:16-nar}); el FFL coherente con compuerta AND filtra pulsos breves de la señal (\cref{fig:16-ffl}); el SIM coordina la expresión de un grupo de genes; dos represores mutuos forman una memoria biestable (\cref{fig:16-toggle}); tres represores en ciclo producen oscilaciones (\cref{fig:16-repr}). Los tres primeros son motivos naturales; los dos últimos, circuitos sintéticos construidos en \especie{E. coli}.}
\label{fig:16-motivos}
\end{figure}

\subsection{El feed-forward loop coherente: un filtro de persistencia}

En un FFL, un regulador $X$ controla a otro regulador $Y$, y ambos controlan conjuntamente a un gen diana $Z$. Como cada una de las tres aristas puede ser activadora o represora, hay ocho tipos. \textcite{c16-mangan2003} analizaron teóricamente los ocho y mostraron que los cuatro \emph{coherentes} (el signo del camino directo coincide con el del indirecto) actúan como \emph{retardos sensibles al signo}, mientras que los cuatro \emph{incoherentes} actúan como \emph{aceleradores sensibles al signo}; también observaron que algunos tipos son mucho más frecuentes que otros en las bases de datos de regulación.

Veamos el más común, el coherente tipo 1 con compuerta AND. Con una señal $S_x$ que activa a $X$ (la tomamos como escalón $X^*\in\{0,1\}$),
\begin{equation}\label{eq:16-ffl}
\frac{dY}{dt}=\beta_y\,f(X^*,K_{xy})-\alpha_y Y,\qquad
\frac{dZ}{dt}=\beta_z\,f(X^*,K_{xz})\,f(Y,K_{yz})-\alpha_z Z,
\end{equation}
\begin{simbolos}
\simbolo{$X^*$}{Forma activa de $X$ (tras recibir su señal).}
\simbolo{$f(\cdot,K)$}{Función de Hill activadora normalizada, $f(u,K)=u^n/(K^n+u^n)$.}
\simbolo{$K_{xy},\ K_{xz},\ K_{yz}$}{Umbrales de activación de cada arista.}
\end{simbolos}
El producto $f(X^*)f(Y)$ implementa la compuerta AND: $Z$ solo se produce si $X^*$ \emph{y} $Y$ están presentes. Cuando la señal aparece, $X^*$ se activa de inmediato, pero $Y$ tiene que acumularse hasta superar $K_{yz}$; solo entonces empieza a producirse $Z$. Si la señal desaparece antes, $Z$ nunca se enciende. Cuando la señal se apaga, en cambio, basta con que $X^*$ desaparezca para que la compuerta AND se cierre: no hay retardo. El circuito ignora los pulsos breves y responde a los persistentes (\cref{fig:16-ffl}), lo que es útil para no activar un programa costoso, como el de utilización de un azúcar, ante una fluctuación pasajera.

\begin{figure}[t]
\centering
\begin{tikzpicture}
\begin{groupplot}[group style={group size=2 by 1, horizontal sep=1.3cm},
   curso, every axis plot/.append style={thick}, width=0.53\linewidth, height=5.2cm, xmin=0, xmax=8, ymin=-0.02, ymax=1.12,
   xlabel={tiempo ($1/\alpha$)}]
\nextgroupplot[title={pulso breve (0,5)}, ylabel={nivel relativo}, legend pos=north east]
  \addplot[gris!70, const plot, fill=gris!12, draw=gris!70] table[x=t,y=Xc] {figuras/cap16/ffl.dat} \closedcycle;
  \addlegendentry{señal $X^*$}
  \addplot[azul] table[x=t,y=Yc] {figuras/cap16/ffl.dat}; \addlegendentry{$Y$}
  \addplot[naranja] table[x=t,y=Zc] {figuras/cap16/ffl.dat}; \addlegendentry{$Z$}
  \addplot[tinta2, densely dotted, thick, forget plot] coordinates {(0,0.5) (8,0.5)};
\nextgroupplot[title={pulso largo (4)}]
  \addplot[gris!70, const plot, fill=gris!12, draw=gris!70] table[x=t,y=Xl] {figuras/cap16/ffl.dat} \closedcycle;
  \addplot[azul] table[x=t,y=Yl] {figuras/cap16/ffl.dat};
  \addplot[naranja] table[x=t,y=Zl] {figuras/cap16/ffl.dat};
  \addplot[tinta, dashed, thick] table[x=t,y=Zs] {figuras/cap16/ffl.dat};
  \addplot[tinta2, densely dotted, thick] coordinates {(0,0.5) (8,0.5)};
  \node[etiqueta, anchor=west, align=left] at (axis cs:5.2,0.85) {discontinua:\\$Z$ sin FFL\\($=Y$)};
\end{groupplot}
\end{tikzpicture}
\caption[El FFL coherente filtra pulsos breves]{Simulación del FFL coherente tipo 1 con compuerta AND (\cref{eq:16-ffl}, $n=4$, $K_{yz}=0{,}5$, $\alpha_y=\alpha_z=1$). \textbf{Izquierda}: ante un pulso de duración $0{,}5$, $Y$ no llega al umbral $K_{yz}$ (línea punteada) y $Z$ apenas alcanza un 3\,\% de su máximo: el pulso se filtra. \textbf{Derecha}: ante un pulso de duración 4, $Z$ se enciende con retraso (tarda $1{,}46$ en llegar a la mitad de su máximo, frente a $0{,}69$ con regulación simple, línea discontinua, que coincide con la curva de $Y$ porque $Y$ tiene justamente regulación simple) pero se apaga sin retraso ($0{,}70$) cuando la señal desaparece: un retardo sensible al signo.}
\label{fig:16-ffl}
\end{figure}

\subsection{El interruptor biestable: nulclinas y estabilidad}

\textcite{c16-gardner2000} construyeron en \especie{E. coli} un circuito sintético con dos promotores reprimibles dispuestos en inhibición mutua: el producto de cada gen reprime al otro. El resultado fue un \emph{interruptor genético}: una red biestable que se puede conmutar entre sus dos estados con un pulso transitorio de un inductor químico o de temperatura, y que recuerda su estado cuando el pulso termina. Un modelo adimensional del circuito es
\begin{equation}\label{eq:16-toggle}
\frac{du}{dt}=\frac{\alpha_1}{1+v^{\beta}}-u,\qquad
\frac{dv}{dt}=\frac{\alpha_2}{1+u^{\gamma}}-v,
\end{equation}
\begin{simbolos}
\simbolo{$u,\ v$}{Concentraciones de los dos represores, en unidades de sus constantes de represión.}
\simbolo{$\alpha_1,\ \alpha_2$}{Tasas de síntesis efectivas (fuerza de cada promotor).}
\simbolo{$\beta,\ \gamma$}{Cooperatividad de la represión de cada promotor.}
\end{simbolos}
con el tiempo medido en unidades de $1/\alpha$ (la tasa de eliminación, que suponemos igual para ambos).

\paragraph{Nulclinas y puntos fijos} Un sistema de dos variables se entiende mejor en su \emph{plano de fases}, donde cada punto $(u,v)$ es un estado y la \cref{eq:16-toggle} asigna a cada uno una velocidad. La \emph{nulclina} de $u$ es la curva donde $du/dt=0$, es decir, $u=\alpha_1/(1+v^\beta)$; la de $v$ es $v=\alpha_2/(1+u^\gamma)$. Donde se cortan, ambas derivadas se anulan: son los \emph{puntos fijos}. Para decidir si un punto fijo es estable, linealizamos.

\begin{teorema}{Estabilidad lineal en dos dimensiones}{16-estabilidad}
Sea $(u^*,v^*)$ un punto fijo de $\dot{\mathbf{x}}=\mathbf{F}(\mathbf{x})$ y $J$ su matriz jacobiana, $J_{ij}=\partial F_i/\partial x_j$ evaluada en él. Las pequeñas perturbaciones evolucionan como $\delta\mathbf{x}(t)=\sum_k c_k\,\mathbf{e}_k\,e^{\lambda_k t}$, donde $\lambda_k$ y $\mathbf{e}_k$ son los autovalores y autovectores de $J$. El punto fijo es estable si todos los autovalores tienen parte real negativa; en dos dimensiones, esto equivale a
\begin{equation}\label{eq:16-trdet}
\operatorname{tr}J<0\quad\text{y}\quad\det J>0,\qquad
\lambda_{1,2}=\frac{\operatorname{tr}J\pm\sqrt{(\operatorname{tr}J)^2-4\det J}}{2}.
\end{equation}
Si $\det J<0$, los autovalores son reales y de signo opuesto: el punto fijo es una \emph{silla}.
\end{teorema}
\begin{simbolos}
\simbolo{$J$}{Jacobiano: cómo responde cada tasa de cambio a pequeñas variaciones de cada variable.}
\simbolo{$\lambda_k$}{Autovalores de $J$; su parte real da la tasa de crecimiento o decaimiento de cada modo.}
\simbolo{$\operatorname{tr}J,\ \det J$}{Traza ($\lambda_1+\lambda_2$) y determinante ($\lambda_1\lambda_2$) de $J$.}
\end{simbolos}

Para el interruptor, derivando la \cref{eq:16-toggle},
\[
J=\begin{pmatrix}-1 & -\dfrac{\alpha_1\beta v^{\beta-1}}{(1+v^\beta)^2}\\[10pt] -\dfrac{\alpha_2\gamma u^{\gamma-1}}{(1+u^\gamma)^2} & -1\end{pmatrix}
\equiv\begin{pmatrix}-1 & -a\\ -b & -1\end{pmatrix},
\qquad \lambda=-1\pm\sqrt{ab}.
\]
La traza siempre vale $-2$; la estabilidad depende solo de si $ab<1$ (estable) o $ab>1$ (silla). En palabras: un punto fijo es inestable cuando la represión mutua es tan sensible que una pequeña ventaja de un represor se amplifica en el otro más de lo que la eliminación puede corregir.

\begin{ejemplo}{Biestabilidad con $\alpha_1=\alpha_2=10$ y $\beta=\gamma=2$}{16-toggle}
Sustituyendo $v=10/(1+u^2)$ en la nulclina de $u$ y resolviendo numéricamente, aparecen tres puntos fijos:
\begin{center}
\small
\begin{tabular}{@{}lcccl@{}}
\toprule
punto fijo $(u^*,v^*)$ & $a$ & $b$ & $\lambda_{1,2}$ & tipo\\
\midrule
$(0{,}101;\ 9{,}899)$ & $0{,}020$ & $1{,}980$ & $-0{,}8;\ -1{,}2$ & estable (``V encendido'')\\
$(2;\ 2)$ & $1{,}6$ & $1{,}6$ & $+0{,}6;\ -2{,}6$ & silla\\
$(9{,}899;\ 0{,}101)$ & $1{,}980$ & $0{,}020$ & $-0{,}8;\ -1{,}2$ & estable (``U encendido'')\\
\bottomrule
\end{tabular}
\end{center}
La silla está en la diagonal y su autovector estable, a lo largo de ella, forma la \emph{separatriz}: toda trayectoria que empieza por encima termina en el estado ``V encendido'', y toda la que empieza por debajo, en ``U encendido''. En cambio, sin cooperatividad ($\beta=\gamma=1$) las nulclinas se cortan en un único punto, $u^*=v^*=2{,}70$, que es estable: no hay memoria.

¿Cuánta síntesis hace falta? En el caso simétrico ($\alpha_1=\alpha_2=\alpha$, $\beta=\gamma=2$) el punto diagonal cumple $u(1+u^2)=\alpha$ y tiene $a=b=2\alpha u/(1+u^2)^2=2u^2/(1+u^2)$. La condición de silla $a>1$ equivale a $u>1$, es decir, $\alpha>2$. Con promotores más débiles que ese umbral, el interruptor no funciona.
\end{ejemplo}

\begin{figure}[t]
\centering
\begin{tikzpicture}
\begin{axis}[curso, every axis plot/.append style={thick}, width=0.72\linewidth, height=0.72\linewidth, xmin=0, xmax=11, ymin=0, ymax=11,
   xlabel={represor $u$}, ylabel={represor $v$}, grid=none, legend pos=north east,
   legend style={fill=white}]
  \addplot[gris!80, -{Stealth[length=1.3mm]}, thin, forget plot,
     quiver={u=\thisrow{du}, v=\thisrow{dv}}] table[x=u,y=v] {figuras/cap16/toggle_q.dat};
  \addplot[azul, very thick] table[x=u,y=v] {figuras/cap16/toggle_nu.dat};
  \addlegendentry{nulclina $\dot u=0$}
  \addplot[naranja, very thick] table[x=u,y=v] {figuras/cap16/toggle_nv.dat};
  \addlegendentry{nulclina $\dot v=0$}
  \foreach \i in {0,...,7}{
    \addplot[violeta, thick, forget plot] table[x=u,y=v] {figuras/cap16/toggle_tr\i.dat};}
  \addplot[violeta, thick] coordinates {(-1,-1)}; \addlegendentry{trayectorias}
  \addplot[only marks, mark=*, mark size=3.2pt, tinta] coordinates {(0.101,9.899) (9.899,0.101)};
  \addlegendentry{estables}
  \addplot[only marks, mark=*, mark size=3.2pt, fill=white, draw=tinta, thick] coordinates {(2,2)};
  \addlegendentry{silla}
  \addplot[tinta2, dashed, thin, forget plot] coordinates {(0,0) (11,11)};
\end{axis}
\end{tikzpicture}
\caption[Plano de fases del interruptor biestable]{Plano de fases del interruptor genético (\cref{eq:16-toggle}, $\alpha_1=\alpha_2=10$, $\beta=\gamma=2$). Las flechas grises indican la dirección del flujo; las curvas azul y naranja son las nulclinas, que se cortan en tres puntos fijos. Las trayectorias (violeta), integradas desde ocho estados iniciales, terminan en uno de los dos estados estables según el lado de la diagonal discontinua (la separatriz) en el que empiezan. Un pulso de inductor que empuje al sistema al otro lado de la separatriz conmuta el estado, y el sistema lo recuerda al retirarse el pulso.}
\label{fig:16-toggle}
\end{figure}

\begin{codigo}[interruptor.py]
import numpy as np
from scipy.integrate import solve_ivp

a1 = a2 = 10.0
n = 2

def toggle(t, y):
    u, v = y
    return [a1 / (1 + v**n) - u, a2 / (1 + u**n) - v]

for y0 in [(0.3, 2.0), (2.0, 0.3), (4.0, 3.9)]:
    s = solve_ivp(toggle, (0, 30), y0, rtol=1e-8)
    print(y0, "->", s.y[:, -1].round(2))
# (0.3, 2.0) -> [0.1 9.9]   (2.0, 0.3) -> [9.9 0.1]
# (4.0, 3.9) -> [9.9 0.1]   (empieza bajo la diagonal)
\end{codigo}

\subsection{El repressilator: oscilaciones a partir de retroalimentación negativa}

En el mismo número de \emph{Nature} en que apareció el interruptor, \textcite{c16-elowitz2000} presentaron un oscilador sintético, el \ingles{repressilator}: tres represores (LacI, TetR y cI del fago $\lambda$) dispuestos en un ciclo en el que cada uno reprime al siguiente. El circuito inducía periódicamente la síntesis de una proteína fluorescente verde en células individuales, con periodos típicos de horas, más largos que el ciclo de división, de modo que el estado del oscilador se transmitía de una generación a la siguiente; las oscilaciones eran ruidosas, posiblemente por fluctuaciones estocásticas de sus componentes. El modelo adimensional usual distingue el ARNm $m_i$ y la proteína $p_i$ de cada gen:
\begin{equation}\label{eq:16-repr}
\frac{dm_i}{dt}=-m_i+\frac{\alpha}{1+p_{j}^{\,n}}+\alpha_0,\qquad
\frac{dp_i}{dt}=-\beta\,(p_i-m_i),\qquad (i,j)\in\{(1,3),(2,1),(3,2)\}.
\end{equation}
\begin{simbolos}
\simbolo{$m_i,\ p_i$}{ARNm y proteína del gen $i$ (adimensionales).}
\simbolo{$\alpha,\ \alpha_0$}{Transcripción máxima y basal (``fuga'' del promotor reprimido).}
\simbolo{$\beta$}{Cociente entre la tasa de eliminación de la proteína y la del ARNm.}
\simbolo{$n$}{Coeficiente de Hill de la represión.}
\end{simbolos}
El tiempo se mide en unidades del tiempo de vida del ARNm.

\paragraph{¿Por qué oscila?} Un ciclo de tres represiones es, en conjunto, una retroalimentación \emph{negativa} (tres signos negativos dan uno negativo) con retraso, porque la señal tarda en recorrer las seis etapas (tres transcripciones, tres traducciones). Si la retroalimentación es lo bastante fuerte, el sistema se ``pasa de largo'' una y otra vez en lugar de asentarse. Hagamos el análisis de estabilidad. Por simetría existe un punto fijo con $m_i=p_i=p^*$, donde $p^*=\alpha/(1+p^{*n})+\alpha_0$. Llamemos $X=-f'(p^*)=\alpha n p^{*\,n-1}/(1+p^{*n})^2>0$ a la sensibilidad de la represión. Linealizando la \cref{eq:16-repr}, $(\lambda+1)\,\delta m_i=-X\,\delta p_j$ y $(\lambda+\beta)\,\delta p_i=\beta\,\delta m_i$; combinándolas y recorriendo el ciclo,
\[
\left[\frac{(\lambda+1)(\lambda+\beta)}{\beta}\right]^3=-X^3
\;\Longrightarrow\;
\frac{(\lambda+1)(\lambda+\beta)}{\beta}\in\left\{-X,\;X e^{i\pi/3},\;X e^{-i\pi/3}\right\}.
\]
La raíz real $-X$ solo da autovalores con parte real negativa. Las complejas pueden cruzar el eje imaginario: sustituyendo $\lambda=i\omega$ en $(\lambda+1)(\lambda+\beta)=\beta X e^{i\pi/3}$ y separando parte real e imaginaria se obtienen $\beta-\omega^2=\beta X/2$ y $\omega(1+\beta)=\beta X\sqrt3/2$. Eliminando $\omega$ llegamos al criterio siguiente.

\begin{teorema}{Condición de oscilación del repressilator}{16-repr}
El punto fijo simétrico de la \cref{eq:16-repr} es inestable, y el sistema oscila, si y solo si
\begin{equation}\label{eq:16-hopf}
\frac{(\beta+1)^2}{\beta}<\frac{3X^2}{4-2X}\qquad\text{o bien}\qquad X\ge2.
\end{equation}
\end{teorema}
\begin{simbolos}
\simbolo{$X$}{Pendiente (cambiada de signo) de la función de represión en el punto fijo.}
\simbolo{$(\beta+1)^2/\beta$}{Término mínimo (igual a 4) cuando $\beta=1$: oscilar es más fácil si el ARNm y la proteína viven lo mismo.}
\end{simbolos}
Hemos comprobado la \cref{eq:16-hopf} contra los autovalores del jacobiano de $6\times6$ en \num{1600} combinaciones de $X$ y $\beta$, sin ninguna discrepancia. Dos consecuencias saltan a la vista. Primera: sin cooperatividad ($n=1$, $\alpha_0=0$) se tiene $p^*(1+p^*)=\alpha$ y $X=p^*/(1+p^*)<1$, de modo que el lado derecho es menor que $3/2$ y nunca supera al izquierdo, que vale al menos 4: \emph{un repressilator con $n=1$ no puede oscilar}. Segunda: las oscilaciones exigen transcripción fuerte (con $n=2$, $\alpha$ debe superar $4{,}24$ incluso en el caso óptimo $\beta=1$) y vidas medias de ARNm y proteína comparables.

\begin{figure}[t]
\centering
\begin{tikzpicture}
\begin{groupplot}[group style={group size=2 by 1, horizontal sep=1.6cm},
   curso, height=5.2cm]
\nextgroupplot[width=0.6\linewidth, xmin=0, xmax=60, ymin=0, ymax=65,
   xlabel={tiempo (vidas del ARNm)}, ylabel={proteína $p_i$},
   legend style={at={(0.5,-0.32)},anchor=north}, legend columns=3, title={oscilaciones}]
  \addplot[azul] table[x=t,y=p1] {figuras/cap16/repressilator.dat}; \addlegendentry{LacI}
  \addplot[naranja] table[x=t,y=p2] {figuras/cap16/repressilator.dat}; \addlegendentry{TetR}
  \addplot[aqua] table[x=t,y=p3] {figuras/cap16/repressilator.dat}; \addlegendentry{cI}
\nextgroupplot[width=0.44\linewidth, xmode=log, ymode=log, xmin=0.1, xmax=1000, ymin=1, ymax=1e5,
   xlabel={$\beta$}, ylabel={$\alpha$}, title={diagrama de estabilidad}, grid=none,
   xtick={0.1,1,10,100,1000}, ytick={1,100,10000}]
  \addplot[name path=fr, tinta2, thick] table[x=beta,y=alfa] {figuras/cap16/repr_frontera.dat};
  \path[name path=top] (axis cs:0.1,1e5) -- (axis cs:1000,1e5);
  \addplot[azulpalido] fill between[of=fr and top];
  \addplot[only marks, mark=*, mark size=2.5pt, naranja] coordinates {(5,216)};
  \node[etiqueta, text=azulprofundo] at (axis cs:3,8000) {oscila};
  \node[etiqueta] at (axis cs:10,4) {estable};
\end{groupplot}
\end{tikzpicture}
\caption[Oscilaciones del repressilator]{\textbf{Izquierda}: simulación de la \cref{eq:16-repr} con parámetros ilustrativos $\alpha=216$, $\alpha_0=0{,}216$, $\beta=5$, $n=2$. Tras un transitorio, las tres proteínas oscilan con un periodo de $8{,}7$ unidades de tiempo, desfasadas un tercio de ciclo: cuando LacI cae, TetR se libera, lo que reprime a cI, lo que a su vez libera a LacI. \textbf{Derecha}: frontera de estabilidad en el plano $(\beta,\alpha)$ para $n=2$ según la \cref{eq:16-hopf}; en la región azul el punto fijo es inestable y el sistema oscila. El punto naranja marca los parámetros de la izquierda; allí el jacobiano tiene un par de autovalores $0{,}208\pm1{,}266i$ con parte real positiva.}
\label{fig:16-repr}
\end{figure}

\subsection{Ruido: cuando las moléculas son pocas}

Las ecuaciones diferenciales describen concentraciones continuas y deterministas. Pero en una bacteria muchos factores de transcripción están presentes en decenas de copias, y sus ARNm en unas pocas. A esas escalas, cada evento de síntesis o degradación es un suceso aleatorio, y dos células genéticamente idénticas en el mismo ambiente difieren. \textcite{c16-elowitz2002} construyeron cepas de \especie{E. coli} con dos reporteros fluorescentes idénticos controlados por el mismo promotor, lo que les permitió separar el ruido \emph{intrínseco} (inherente a la expresión de cada gen, que hace diferir a los dos reporteros dentro de una célula) del \emph{extrínseco} (fluctuaciones de otros componentes celulares, que los mueve juntos). Ambos contribuían sustancialmente, y los bajos números de copia imponían un límite a la precisión de la regulación.

\paragraph{La ecuación maestra} El modelo más sencillo es un nacimiento y muerte: la proteína se produce a tasa $k$ y cada molécula se degrada a tasa $\gamma$. La probabilidad $P_x(t)$ de tener $x$ moléculas cumple
\begin{equation}\label{eq:16-cme}
\frac{dP_x}{dt}=k\,P_{x-1}+\gamma\,(x+1)\,P_{x+1}-(k+\gamma x)\,P_x .
\end{equation}
\begin{simbolos}
\simbolo{$P_x(t)$}{Probabilidad de que haya exactamente $x$ moléculas en el tiempo $t$.}
\simbolo{$k$}{Tasa de producción (moléculas por unidad de tiempo).}
\simbolo{$\gamma$}{Tasa de degradación por molécula.}
\end{simbolos}
En el estado estacionario, el flujo de $x-1$ a $x$ equilibra el de $x$ a $x-1$: $kP_{x-1}=\gamma x P_x$, de donde $P_x=(\mu/x)P_{x-1}$ con $\mu=k/\gamma$, y por inducción $P_x=e^{-\mu}\mu^x/x!$. La distribución es de Poisson: su media coincide con la solución de la ODE, y su coeficiente de variación es $\mathrm{CV}=\sigma/\mu=1/\sqrt{\mu}$. Con 10 moléculas de media, las fluctuaciones típicas son del 32\,\%; con \num{10000}, del 1\,\%.

\paragraph{El algoritmo de Gillespie} Para redes más complejas la ecuación maestra no se puede resolver, pero sí se pueden generar trayectorias exactas. \textcite{c16-gillespie1977} observó que, si el sistema está en un estado con \emph{propensiones} $a_1,\dots,a_R$ (la probabilidad por unidad de tiempo de que ocurra cada reacción), el tiempo hasta la siguiente reacción es exponencial con tasa $a_0=\sum_r a_r$, y la reacción que ocurre es la $r$-ésima con probabilidad $a_r/a_0$:
\begin{equation}\label{eq:16-gillespie}
\tau=\frac{1}{a_0}\ln\frac{1}{u_1},\qquad
r=\min\Big\{r':\ \sum_{s\le r'}a_s> u_2\,a_0\Big\},\qquad u_1,u_2\sim\mathcal{U}(0,1).
\end{equation}
\begin{simbolos}
\simbolo{$a_r$}{Propensión de la reacción $r$ en el estado actual (p.\,ej., $k$ para la síntesis y $\gamma x$ para la degradación).}
\simbolo{$\tau$}{Tiempo de espera hasta la próxima reacción.}
\simbolo{$u_1,u_2$}{Números aleatorios uniformes independientes.}
\end{simbolos}
El algoritmo es exacto en el sentido de que las trayectorias generadas tienen exactamente la distribución que describe la ecuación maestra.

\begin{codigo}[gillespie.py]
import numpy as np

def gillespie_nm(k, g, x0, T, rng):
    """Nacimiento (tasa k) y muerte (tasa g*x) de una proteína."""
    t, x = 0.0, x0
    ts, xs = [t], [x]
    while t < T:
        a1, a2 = k, g * x                  # propensiones
        a0 = a1 + a2
        t += rng.exponential(1 / a0)       # ¿cuándo?
        x += 1 if rng.random() * a0 < a1 else -1   # ¿cuál?
        ts.append(t)
        xs.append(x)
    return np.array(ts), np.array(xs)

ts, xs = gillespie_nm(1.0, 0.1, 10, 20000, np.random.default_rng(7))
\end{codigo}

\begin{figure}[t]
\centering
\begin{tikzpicture}
\begin{groupplot}[group style={group size=2 by 1, horizontal sep=1.5cm},
   curso, height=5.2cm]
\nextgroupplot[width=0.6\linewidth, xmin=0, xmax=100, ymin=0, ymax=24,
   xlabel={tiempo (min)}, ylabel={moléculas $x$}, title={trayectorias}, legend pos=south east,
   legend columns=2]
  \addplot[azul!70, thin, const plot] table[x=t,y=x] {figuras/cap16/ssa0.dat}; \addlegendentry{Gillespie}
  \addplot[aqua!80!black, thin, const plot, forget plot] table[x=t,y=x] {figuras/cap16/ssa1.dat};
  \addplot[violeta!70, thin, const plot, forget plot] table[x=t,y=x] {figuras/cap16/ssa2.dat};
  \addplot[naranja, very thick] table[x=t,y=x] {figuras/cap16/ssa_ode.dat}; \addlegendentry{ODE}
\nextgroupplot[width=0.44\linewidth, xmin=-0.5, xmax=25, ymin=0, ymax=0.14,
   xlabel={$x$}, ylabel={probabilidad}, title={distribución estacionaria},
   yticklabel style={/pgf/number format/fixed, /pgf/number format/precision=2}]
  \addplot[ybar, bar width=2.4pt, fill=azul!55, draw=none] table[x=x,y=ssa] {figuras/cap16/ssa_hist.dat};
  \addplot[only marks, mark=o, mark size=1.6pt, naranja, thick] table[x=x,y=pois] {figuras/cap16/ssa_hist.dat};
\end{groupplot}
\end{tikzpicture}
\caption[Simulación estocástica frente a determinista]{\textbf{Izquierda}: tres trayectorias exactas del algoritmo de Gillespie (\cref{eq:16-gillespie}) para una proteína con síntesis $k=1\ \text{min}^{-1}$ y degradación $\gamma=0{,}1\ \text{min}^{-1}$, partiendo de cero, frente a la solución de la ODE, $x(t)=10(1-e^{-0{,}1t})$ (naranja). La ODE describe bien la media, pero no las fluctuaciones de célula a célula. \textbf{Derecha}: distribución del número de moléculas en una simulación larga (\num{39960} reacciones, barras) frente a la distribución de Poisson de media 10 que predice la \cref{eq:16-cme} (círculos). La simulación da media $10{,}05$, varianza $10{,}19$ y $\mathrm{CV}=0{,}318$, frente al $1/\sqrt{10}=0{,}316$ teórico.}
\label{fig:16-ssa}
\end{figure}

\begin{profundizacion}[Ruido que decide]
El ruido no es solo una molestia. En un sistema biestable como el de la \cref{fig:16-toggle}, las fluctuaciones pueden empujar a una célula a través de la separatriz y conmutar su estado sin ninguna señal externa. Así, una población clonal puede dividirse en dos subpoblaciones con fenotipos distintos, una estrategia de ``apuesta diversificada'' ante ambientes impredecibles. Simular estos fenómenos exige el algoritmo de Gillespie (o sus aproximaciones más rápidas, como el \ingles{tau-leaping}) aplicado al circuito completo, porque la ODE, que solo ve la media, predice que cada célula se queda para siempre en el estado donde empezó.
\end{profundizacion}

\subsection{Inferir la red a partir de los datos}

Hasta aquí supusimos conocido el circuito. En la práctica, a menudo tenemos lo contrario: una matriz de expresión (genes por muestras, como las de RNA-seq del \cref{cap:11}) y la pregunta de qué gen regula a cuál. La idea más ingenua es unir los genes muy correlacionados, pero la correlación no distingue un efecto directo ($X\to Z$) de uno indirecto ($X\to Y\to Z$) ni de una causa común.

\textcite{c16-huynhthu2010} propusieron GENIE3, que descompone el problema de inferir una red de $p$ genes en $p$ problemas de regresión. Para cada gen diana $j$, se predice su expresión a partir de la de todos los reguladores candidatos con un conjunto de árboles (\ingles{Random Forests} o \ingles{Extra-Trees}), y la importancia de cada regulador en esa predicción se toma como evidencia de una arista:
\begin{equation}\label{eq:16-genie3}
x_j = f_j\big(\mathbf{x}_{-j}\big)+\varepsilon_j,\qquad
w_{ij}=\mathrm{Imp}_i\big(\hat f_j\big),
\end{equation}
\begin{simbolos}
\simbolo{$x_j$}{Expresión del gen diana $j$ en cada muestra.}
\simbolo{$\mathbf{x}_{-j}$}{Expresión de los reguladores candidatos (excluido $j$).}
\simbolo{$\hat f_j$}{Conjunto de árboles ajustado para predecir $x_j$.}
\simbolo{$w_{ij}$}{Importancia del regulador $i$ (reducción total de varianza en los nodos de los árboles que lo usan).}
\end{simbolos}
Las importancias de todas las dianas se ordenan en una lista global de aristas candidatas. El método no supone ninguna forma particular de regulación, capta interacciones combinatorias y no lineales, produce redes dirigidas y escala bien; fue el mejor en el desafío DREAM4 \ingles{In Silico Multifactorial}.

\begin{codigo}[genie3.py]
import numpy as np
from sklearn.ensemble import RandomForestRegressor

def genie3(X, reguladores, n_arboles=300):
    """X: muestras x genes (estandarizada). Devuelve W[i, j]."""
    n_genes = X.shape[1]
    W = np.zeros((n_genes, n_genes))
    for j in range(n_genes):
        entradas = [i for i in reguladores if i != j]
        rf = RandomForestRegressor(n_estimators=n_arboles,
                                   max_features="sqrt")
        rf.fit(X[:, entradas], X[:, j])
        W[entradas, j] = rf.feature_importances_
    return W
\end{codigo}

¿Cómo se evalúa un método de inferencia? Con una red de referencia, se recorre la lista ordenada de aristas y se calcula, para cada umbral, la \emph{precisión} (fracción de aristas predichas que son reales) y la \emph{exhaustividad} (fracción de aristas reales recuperadas); el área bajo esa curva (AUPR) resume el desempeño. La \cref{fig:16-pr} muestra un experimento controlado: generamos redes aleatorias de 20 genes con 6 reguladores, simulamos 300 estados estacionarios bajo perturbaciones aleatorias y comparamos GENIE3 con la correlación absoluta. En promedio sobre cinco redes, GENIE3 obtiene AUPR $=0{,}54\pm0{,}06$ y la correlación $0{,}53\pm0{,}09$, ambos el doble de lo esperable al azar ($0{,}27$). La ventaja de GENIE3 es pequeña y no aparece en todas las redes.

\begin{figure}[t]
\centering
\begin{tikzpicture}
\begin{axis}[curso, every axis plot/.append style={thick}, width=0.78\linewidth, height=5.4cm, xmin=0, xmax=1, ymin=0, ymax=1.05,
   xlabel={exhaustividad (\ingles{recall})}, ylabel={precisión}, legend pos=north east]
  \addplot[azul] table[x=r,y=p] {figuras/cap16/pr_genie3.dat};
  \addlegendentry{GENIE3 (AUPR $0{,}60$)}
  \addplot[naranja] table[x=r,y=p] {figuras/cap16/pr_cor.dat};
  \addlegendentry{$|$correlación$|$ (AUPR $0{,}57$)}
  \addplot[tinta2, dashed, thick] coordinates {(0,0.281) (1,0.281)};
  \addlegendentry{azar ($0{,}28$)}
\end{axis}
\end{tikzpicture}
\caption[Curvas de precisión-exhaustividad]{Evaluación de la inferencia en una red sintética de 20 genes (6 reguladores, 32 aristas verdaderas entre 114 candidatas) a partir de 300 perfiles de expresión simulados. Ambos métodos recuperan las primeras aristas con precisión alta, pero la precisión cae rápidamente: a mitad de exhaustividad, aproximadamente la mitad de las predicciones son falsas. La línea discontinua es la precisión de una lista ordenada al azar. Simulación en \archivo{figuras/cap16/generar.py}.}
\label{fig:16-pr}
\end{figure}

Un resultado tan modesto no es excepcional en la inferencia de redes. En el proyecto DREAM5, \textcite{c16-marbach2012} evaluaron a ciegas más de 30 métodos de inferencia sobre datos de \especie{E. coli}, \especie{Staphylococcus aureus}, \especie{S. cerevisiae} y datos simulados, y observaron que ningún método era el mejor en todos los conjuntos de datos. En cambio, \emph{integrar} las predicciones de muchos métodos (la ``sabiduría de las multitudes'') daba un desempeño alto y robusto, y con ello construyeron redes de alta confianza para \especie{E. coli} y \especie{S. aureus}, de unas \num{1700} interacciones transcripcionales cada una, con una precisión de alrededor del 50\,\%.

\begin{cuidado}[Lo que los datos de expresión no pueden decir]
Los datos observacionales de expresión, por sí solos, rara vez identifican la causalidad: dos genes corregulados por un tercero no medido parecen regularse entre sí, y un factor de transcripción cuyo ARNm apenas cambia (porque se regula por fosforilación) es invisible. Las perturbaciones (\ingles{knockouts}, CRISPRi), las series temporales y los datos de unión al ADN (ChIP-seq, motivos del \cref{cap:04}) restringen enormemente el espacio de redes compatibles. Una red inferida es un conjunto de hipótesis ordenadas por plausibilidad, no un mapa.
\end{cuidado}

\begin{ideaclave}
La topología dice qué \emph{puede} hacer un circuito; los parámetros (cooperatividad, fuerzas de promotor, vidas medias) deciden qué \emph{hace}. Un mismo ciclo de represores es estable u oscila según $n$, $\alpha$ y $\beta$.
\end{ideaclave}

\begin{resumen}
\begin{itemize}
  \item Un interactoma es un grafo; su matriz de adyacencia codifica grados ($k_i=\sum_jA_{ij}$), caminos ($(A^\ell)_{ij}$) y triángulos ($(A^3)_{ii}/2$), de los que salen el coeficiente de agrupamiento, las distancias y la intermediación.
  \item El doble híbrido mide pares binarios y AP-MS mide complejos (modelos estrella o matriz); cada método tiene su perfil de falsos positivos y negativos. STRING integra canales de evidencia y BioGRID cura interacciones de la literatura.
  \item Las redes biológicas son mundos pequeños muy agrupados (Watts-Strogatz). El enlace preferencial de Barabási-Albert produce $P(k)\propto k^{-3}$, pero una cola ancha en un gráfico log-log no demuestra una ley de potencias; en la red de levadura, los grandes complejos dominan la cola.
  \item La modularidad $Q$ compara las aristas internas de una partición con las esperadas por azar; Louvain la maximiza y las comunidades deben validarse con anotaciones y redes aleatorizadas.
  \item En el modelo $\dot Y=\beta-\alpha Y$, el tiempo de respuesta es $\ln2/\alpha$, independiente de la producción. La autorregulación negativa lo acorta y amortigua el ruido; el FFL coherente con AND filtra pulsos breves.
  \item La represión mutua cooperativa crea biestabilidad (dos puntos estables y una silla, identificados por los autovalores del jacobiano); un ciclo de tres represores oscila si se cumple $(\beta+1)^2/\beta<3X^2/(4-2X)$.
  \item Con pocas moléculas la expresión es estocástica; el algoritmo de Gillespie simula trayectorias exactas y, en nacimiento-muerte, reproduce la distribución de Poisson con $\mathrm{CV}=1/\sqrt{\mu}$.
  \item Inferir redes a partir de la expresión es difícil: GENIE3 convierte el problema en regresiones con árboles, y la integración de muchos métodos es más robusta que cualquiera de ellos por separado.
\end{itemize}
\end{resumen}

\begin{ejercicios}
\ej{1} Para el grafo de la \cref{fig:16-grafo}, calcule $A^2$ y verifique que su diagonal es el vector de grados. ¿Qué cuenta $(A^2)_{ij}$ cuando $i\neq j$?
\ej{1} Una purificación AP-MS recupera un cebo y 11 presas. ¿Cuántas aristas genera el modelo estrella y cuántas el modelo matriz? Si el complejo real tiene forma de anillo, ¿qué fracción de las aristas del modelo matriz son contactos físicos?
\ej{1} Una proteína de \especie{E. coli} estable tiene $t_{1/2}$ igual a un tiempo de generación. ¿Cuál debería ser la vida media de degradación activa para que el tiempo de respuesta sea de un décimo de generación? ¿Cuánto debe aumentar la síntesis para conservar el nivel estacionario?
\ej{2} Con \cmd{networkx}, genere una red ER y otra BA con $n=2000$ y $\langle k\rangle\approx8$. Compare $C$, $L$ y la distribución de grado con la de la \cref{fig:16-grado}. ¿Qué modelo se parece más a la red de levadura y en qué falla cada uno?
\ej{2} Descargue de STRING la red física de su organismo favorito, extraiga la ego-red de radio 2 de una proteína de interés, detecte comunidades con Louvain y compare $Q$ con el de 20 redes aleatorizadas con \py{nx.double_edge_swap}. Pruebe varias semillas: ¿cambian las comunidades?
\ej{2} Demuestre la \cref{eq:16-narest} y compruebe numéricamente, integrando la \cref{eq:16-nar}, cuánto cambia $X_{\text{est}}$ al duplicar $\beta$ con $n=1$, $2$ y $4$.
\ej{2} Simule el FFL de la \cref{eq:16-ffl} con pulsos de duración creciente y grafique el máximo de $Z$ en función de la duración. Estime la duración mínima de pulso que el circuito deja pasar.
\ej{3} Para el interruptor simétrico con $n=3$, encuentre el valor crítico de $\alpha$ a partir del cual aparecen tres puntos fijos, generalizando el razonamiento del \cref{ej:16-toggle}.
\ej{3} Complete la derivación de la \cref{eq:16-hopf}: sustituya $\lambda=i\omega$, elimine $\omega$ y muestre que para $X\ge2$ siempre existe un autovalor con parte real positiva. Verifique el periodo de la \cref{fig:16-repr} con la frecuencia $\omega$ del par de autovalores dominante y explique la discrepancia.
\ej{3} Implemente el algoritmo de Gillespie para el interruptor biestable (cuatro reacciones: síntesis y degradación de $u$ y de $v$, con propensiones $\Omega\alpha/(1+(v/\Omega)^2)$ y $u$, etc.) y mida el tiempo medio de conmutación espontánea en función del tamaño del sistema $\Omega$.
\end{ejercicios}

\referenciascapitulo
